\documentclass[11pt]{amsart}
\usepackage[T1]{fontenc}
\usepackage[utf8]{inputenc}
\usepackage{lmodern}
\usepackage[margin=1.05in,headheight=14pt]{geometry}
\usepackage{amsmath,amssymb,mathtools,mathrsfs}
\usepackage{microtype,enumitem,booktabs,array,longtable}
\usepackage{xcolor}
\definecolor{linkblue}{HTML}{1C4765}
\usepackage[colorlinks=true,linkcolor=linkblue,citecolor=linkblue,urlcolor=linkblue]{hyperref}
\usepackage{bookmark}
\numberwithin{equation}{section}
\allowdisplaybreaks[2]
\setlength{\emergencystretch}{2em}
\setlist{itemsep=2pt,topsep=4pt}
\newtheorem{theorem}{Theorem}[section]
\newtheorem{lemma}[theorem]{Lemma}
\newtheorem{proposition}[theorem]{Proposition}
\newtheorem{corollary}[theorem]{Corollary}
\theoremstyle{definition}
\newtheorem{definition}[theorem]{Definition}
\newtheorem{example}[theorem]{Example}
\theoremstyle{remark}
\newtheorem{remark}[theorem]{Remark}
\newcommand{\C}{\mathbb C}
\newcommand{\R}{\mathbb R}
\newcommand{\N}{\mathbb N}
\newcommand{\Z}{\mathbb Z}
\newcommand{\E}{\mathbb E}
\newcommand{\dd}{\,\mathrm d}
\newcommand{\Id}{\mathrm{Id}}
\newcommand{\repourl}{https://github.com/veridiscoverLab/fermionic-coherent-lean}
\DeclareMathOperator{\Tr}{Tr}
\DeclareMathOperator{\rank}{rank}
\DeclareMathOperator{\Span}{span}
\DeclareMathOperator{\supp}{supp}
\DeclareMathOperator{\Spin}{Spin}
\DeclareMathOperator{\SO}{SO}
\DeclareMathOperator{\GL}{GL}
\DeclareMathOperator{\SL}{SL}
\DeclareMathOperator{\Sp}{Sp}
\DeclareMathOperator{\Ann}{Ann}
\DeclareMathOperator{\rev}{rev}
\DeclareMathOperator{\vol}{vol}
\DeclareMathOperator{\diag}{diag}
\DeclareMathOperator{\Pf}{Pf}
\title[Fermionic coherent states]{Fermionic coherent states:\\convex order and Gaussian rank}
\author{veridiscoverLab}
\date{September 10, 2026}
\subjclass[2020]{81R30, 81V74, 60E15, 14M17}
\keywords{Fermionic coherent states, Husimi distribution, Wehrl entropy,
Gaussian rank, pure spinors, Lean formalization}
\hypersetup{pdftitle={Fermionic coherent states: convex order and Gaussian rank},
pdfauthor={veridiscoverLab},pdfsubject={Convex order, fermionic entropy, and Gaussian decompositions}}
\begin{document}
\begin{abstract}
We study two extremal properties of fermionic coherent states: concentration
of their overlap distributions and length of exact Gaussian decompositions.
For every basic exterior-power representation, the Husimi distribution of
any density operator is dominated in convex order by that of a Slater state.
For each strictly convex test function, equality characterizes pure Slater
states. We prove the corresponding statements for the basic half-spin and
odd-dimensional spin representations, obtain all positive-order
R\'enyi--Wehrl minima, and evaluate the coherent distributions and entropy
constants. The proof uses exact conditioning and a single convex
decomposition shared by all measured directions. For Gaussian rank, a
degree-eighteen relative invariant excludes every three-term decomposition
of two copies of the four-mode magic state. This yields rank multiplicativity
for two arbitrary nonzero four-mode even vectors. The exterior-power
convex-order theorem, its equality cases, and the entropy minimizer theorems
have end-to-end Lean proofs. The formalization of the Gaussian-rank theorem
currently includes the original rank bounds two and four and the full
normal-family polynomial certificate.
\end{abstract}
\maketitle
\section{Introduction}
\label{sec:introduction}

A fermionic state can be compared with simple states in two different ways.
One can measure its overlaps with them, or ask how many are needed to express
it as a sum. The first question concerns a probability distribution; the
second concerns an exact vector identity. This paper treats both questions
using the geometry of Slater and Gaussian states.

For a fixed number of particles, the one-particle space is $\C^n$ and the
$p$-particle space is its exterior power $\Lambda^p\C^n$. A unit vector
$u_1\wedge\cdots\wedge u_p$, with the $u_i$ orthonormal, is called a
\emph{Slater vector}. It describes the occupation of $p$ mutually orthogonal
one-particle modes. Even in a small space, most vectors need not have this
form. For example, in $\Lambda^2\C^4$ the normalized vector
\begin{equation}\label{intro:two-particle}
 \psi=\frac{e_1\wedge e_2+e_3\wedge e_4}{\sqrt2}
\end{equation}
is not a Slater vector: its exterior square is nonzero, whereas the square
of a decomposable two-vector vanishes. Its overlap with $e_1\wedge e_2$ has
squared modulus $1/2$. Rotating the reference Slater vector gives an entire
distribution of such overlaps.

A density operator is a positive semidefinite operator $\rho$ of trace one;
it includes both pure states $|\psi\rangle\langle\psi|$ and mixtures. If
$\Omega_x$ is a unit coherent vector indexed by a point $x$ of the coherent
orbit, its \emph{Husimi function} is
\[
 Q_\rho(x)=\langle\Omega_x,\rho\Omega_x\rangle.
\]
The inner product is conjugate-linear in its first argument. The orbit
carries the probability measure $\nu$ induced by normalized Haar measure
of the acting compact group. Thus $Q_\rho$ is a random variable in $[0,1]$.
For a $D$-dimensional irreducible representation, its mean is $1/D$.
For a unit input vector $\psi$, we abbreviate
$Q_\psi=Q_{|\psi\rangle\langle\psi|}$.

\subsection{Comparing the whole distribution}
Our distribution theorem says that coherent inputs maximize every continuous
convex average:
\begin{equation}\label{intro:convex-order}
 \int \Phi(Q_\rho)\dd\nu
 \leq \int \Phi(Q_{P_\Omega})\dd\nu,
 \qquad P_\Omega=|\Omega\rangle\langle\Omega|.
\end{equation}
Here $\Omega$ is any coherent unit vector, and $\Phi$ is continuous and convex
on $[0,1]$. This relation is called \emph{convex order}. Since all these
random variables have the same mean, it says that the coherent distribution
has the greatest spread in every convex test. It implies a comparison of
variances by taking $\Phi(t)=t^2$, and of Wehrl entropies by taking
$\Phi(t)=t\log t$ with the continuous value zero at $t=0$. For every
strictly convex $\Phi$, equality in \eqref{intro:convex-order} forces the
whole input density to be a pure coherent state.

We prove this statement for all basic exterior powers
(Theorem~\ref{ext:convex-order}), for both basic half-spin representations
(Theorem~\ref{spin:convex-order}), and for the basic odd-dimensional spin
representation (Corollary~\ref{spin:odd}). In the half-spin realization
used below, the coherent vectors are pure fermionic Gaussian vectors,
whose particle number may vary while their parity remains fixed.
Section~\ref{sec:spin} constructs this realization and explains its
relation to odd-dimensional spin.

The entropy problem has a long history. Gnutzmann and \.{Z}yczkowski
introduced R\'enyi--Wehrl localization measures in the generalized coherent
state setting~\cite{gnutzmann2001}. Sugita proved coherent-state optimality
for integer indices greater than one~\cite{sugita2002}, and studied the
resulting moments for many-particle states~\cite{sugita2003}. Lieb and
Solovej proved a convex-function comparison for symmetric $\mathrm{SU}(N)$
coherent states~\cite{liebsolovej2016}. Symmetric powers describe a different
representation series from the exterior powers considered here. Our
comparison treats arbitrary continuous convex functions and all mixed
inputs in the stated fermionic series. In particular, it yields all real
positive-order R\'enyi--Wehrl minima, together with their equality cases.

The main step is a compatibility statement about conditioning. Detecting one
occupied mode turns the original density into a density on a smaller
fermionic space. Its normalization is an occupation probability, determined
by a one-particle matrix or a Majorana covariance matrix. We express this
matrix as a convex combination of coherent matrices using one set of
probabilities, chosen from the diagonal of the same rotated full density.
These probabilities work for every subsequent measured direction. The
conditional coherent states also remain coherent in the smaller space.
The induction therefore uses the same comparison state at every stage.
The equality argument recovers the full density from this common
decomposition; it does not identify a many-particle state merely by its
one-particle matrix.

\subsection{Counting Gaussian summands}
Allowing superpositions of particle numbers leads to a different question.
The \emph{Gaussian rank} $\chi_G(\psi)$ is the least number of nonzero
Gaussian vectors whose sum is $\psi$, with arbitrary complex coefficients
absorbed into the summands. The definition and parity convention appear in
Section~\ref{sec:gaussian-rank}.

On four modes, let $|0000\rangle$ and $|1111\rangle$ denote the vacuum and
the state with all four modes occupied, and put
\[
 M=\frac{|0000\rangle+|1111\rangle}{\sqrt2}.
\]
This state is a basic non-Gaussian resource in matchgate computation. Its
two-copy expansion has four Gaussian terms. A shorter expansion, however,
could use terms coupling the two blocks of four modes. Cudby and Strelchuk
proved a rank-four result under symmetry restrictions and formulated the
unrestricted two-copy problem~\cite{cudby2023}. Gaussian decompositions are
also central to the simulation methods of Dias and Koenig~\cite{diaskoenig2024}.
Related multiplicativity results for Gaussian fidelity and extent concern
different optimization problems~\cite{reardon2024}.
Recent work of Tarabunga studies non-Gaussianity through the eigenvalue
distribution of a two-copy Majorana operator and its behavior under
postselected Gaussian operations~\cite{tarabunga2026}.

We prove that $\chi_G(M^{\otimes2})=4$, and more generally that Gaussian
rank is multiplicative for two arbitrary nonzero four-mode even vectors.
The obstruction is a polynomial built from the full bilinear Gram of the
Clifford bivector action. It is nonzero at the target and vanishes on every
sum of three pure spinors. A common change of coordinates first normalizes
two summands. A polynomial deformation then brings the third summand into
a tractable chart, and a specialization argument recovers all degenerate
cases. The change of coordinates acts on the entire sum: transforming the
summands independently would destroy the equation being tested.

\subsection{Formal verification and organization}
The accompanying Lean development verifies the exterior-power
convex-order theorem, the actual Slater-orbit probability measure, the
strict equality characterization, and the positive-order moment and entropy
minimizer theorems. These are end-to-end proofs with no additional
mathematical hypotheses beyond those stated in the theorems. The standard
logical foundations and the precise boundary of the formalization are
recorded in Section~\ref{sec:formalization}. In particular, the Gaussian-rank
four conclusion still has a paper proof supplemented by verified coordinate
certificates and original-object lemmas.

The sources, full declaration types, and verification evidence are available at
\begin{center}
 \small\url{https://github.com/veridiscoverLab/fermionic-coherent-lean}.
\end{center}
We begin with exterior powers, where the conditioning argument is most
transparent, and then pass to spin representations. After evaluating the
coherent distributions and entropy constants, we turn to Gaussian rank.
The final application section gives quantitative entropy bounds and the
information capacity of the specified coherent measurements.

\section{Exterior-power Husimi convex order}
\label{ext:section}

\subsection{Slater states and the measurement}

Let $E=\C^n$ with its standard Hermitian inner product, linear in the
second variable. For $0\le p\le n$, equip
\[
 \mathcal H_{p,n}=\Lambda^p E,
 \qquad D_{p,n}=\dim\mathcal H_{p,n}=\binom np,
\]
with the inner product for which increasing wedges of an orthonormal
basis are orthonormal. Thus a wedge of $p$ orthonormal vectors has norm
one. The induced unitary representation is
$R_p(U)=\Lambda^p U$, for $U\in U(n)$.

If $W\subset E$ is a $p$-dimensional subspace and $(w_1,\ldots,w_p)$
is an orthonormal basis of $W$, the vector
$\Omega_W=w_1\wedge\cdots\wedge w_p$ is a \emph{unit Slater vector}.
A different orthonormal basis changes it by a scalar of modulus one.
Consequently,
\[
 \rho_W=|\Omega_W\rangle\langle\Omega_W|
\]
is determined by $W$. We call $\rho_W$ a \emph{pure Slater density} and
write $\pi_W$ for the orthogonal projection from $E$ onto $W$; these
operators act on different spaces. The Slater-projector orbit is
$\mathcal O_{p,n}=\{\rho_W:\dim W=p\}$.

Fix a unit Slater reference $\Omega_0$. The pushforward of Haar
probability measure on $U(n)$ under
\[
 U\longmapsto
 |R_p(U)\Omega_0\rangle\langle R_p(U)\Omega_0|
\]
is denoted by $\nu_{p,n}$. Transitivity and Haar invariance make this
measure independent of the reference. Equivalently, it is the invariant
probability measure on the Grassmannian of $p$-planes, expressed through
their Slater projectors.

A \emph{density operator} on $\mathcal H_{p,n}$ is an operator
$\rho\ge0$ with $\Tr\rho=1$. Its Husimi function is
\begin{equation}
 Q_\rho(P)=\Tr(\rho P),\qquad P\in\mathcal O_{p,n}.
 \label{ext:husimi}
\end{equation}
We also write
$Q_\rho(U)=\langle R_p(U)\Omega_0,\rho R_p(U)\Omega_0\rangle$
for the same measurement pulled back to the group. Since
$0\le\rho\le\Id$, its values lie in $[0,1]$.

The normalization is
\begin{equation}
 \int_{\mathcal O_{p,n}}P\,d\nu_{p,n}(P)
 =\frac{\Id_{\mathcal H_{p,n}}}{D_{p,n}},
 \qquad
 \int Q_\rho\,d\nu_{p,n}=\frac1{D_{p,n}}.
 \label{ext:resolution}
\end{equation}
Indeed, the averaged projector commutes with every $R_p(U)$.
Commutation with diagonal one-particle unitaries makes it diagonal in
the occupation basis, because distinct occupied subsets have distinct
characters. Permutation unitaries make all its diagonal entries equal.
Its trace is one, which proves the first identity; pairing with $\rho$
proves the second. For $p=0$ and $p=n$, the sector is one-dimensional,
its unique density is Slater, and $Q_\rho\equiv1$. This also includes
$n=p=0$.

\begin{theorem}[Exterior-power Husimi convex order]
\label{ext:convex-order}
For every $0\le p\le n$, every density operator $\rho$ on
$\mathcal H_{p,n}$, every continuous convex function
$\Phi:[0,1]\to\R$, and every pure Slater density $\rho_W$,
\begin{equation}
 \int\Phi(Q_\rho)\,d\nu_{p,n}
 \le
 \int\Phi(Q_{\rho_W})\,d\nu_{p,n}.
 \label{ext:order}
\end{equation}
The right-hand side is independent of $W$. If $\Phi$ is strictly
convex, equality holds if and only if $\rho$ is a pure Slater density.
\end{theorem}

All the compared distributions have the mean in
\eqref{ext:resolution}. The theorem says that a Slater input maximizes
every continuous convex observation of the measurement value. The
equality statement concerns the whole density operator, including its
off-diagonal entries.

\subsection{A shared convex mixture}

We first record the probabilistic step used here and in the spinor
argument. Its coefficients are fixed before the measurement variable
is sampled.

\begin{lemma}[Shared-mixture comparison]
\label{ext:common-mixture}
Let $(X,\mu)$ be a probability space. Suppose measurable functions
$a,a_1,\ldots,a_L:X\to[0,1]$ satisfy
\[
 a(x)=\sum_{j=1}^L t_j a_j(x),\qquad
 t_j\ge0,\qquad\sum_{j=1}^L t_j=1,
\]
and all $a_j$ have the same distribution $\eta$. For every continuous
convex function $H:[0,1]\to\R$,
\begin{equation}
 \int_X H(a(x))\,d\mu(x)
 \le \int_{[0,1]}H(t)\,d\eta(t).
 \label{ext:shared-jensen}
\end{equation}
If $H$ is strictly convex, equality implies that all the $a_j$ with
$t_j>0$ agree almost everywhere.
\end{lemma}
\begin{proof}
Pointwise Jensen gives $H(a)\le\sum_jt_jH(a_j)$. Integration and
the common distribution give \eqref{ext:shared-jensen}. All functions
being integrated are bounded. If the integrals are equal, the
nonnegative Jensen difference vanishes almost everywhere. The equality
criterion for strict convexity gives the assertion.
\end{proof}

If $R$ is a $[0,1]$-valued random variable, define
\begin{equation}
 H_\Phi(t)=\E\Phi(tR),\qquad 0\le t\le1.
 \label{ext:scale-function}
\end{equation}
Continuity follows by bounded convergence, and convexity follows
pointwise. If $\Phi$ is strictly convex and $\Pr(R>0)>0$, then
$H_\Phi$ is strictly convex: for distinct $s,t$ the strict scalar
inequality holds on the event $R>0$.

\subsection{Contraction and the one-particle density}

For $u\in E$, let $c(u)^*$ denote creation by $u$, namely
$\xi\mapsto u\wedge\xi$, and let $c(u)$ be its adjoint. Thus $c(u)$
is conjugate-linear in $u$. These operators act on the full graded
exterior space $\Lambda^\bullet E$ and change degree by one; contraction
vanishes on degree zero. On a decomposable vector,
\begin{equation}
 c(u)(v_1\wedge\cdots\wedge v_p)
 =\sum_{j=1}^p(-1)^{j-1}\langle u,v_j\rangle
 v_1\wedge\cdots\wedge\widehat v_j\wedge\cdots\wedge v_p.
 \label{ext:contraction}
\end{equation}
The hat indicates the omitted factor. The canonical anticommutation
relations are
\[
 c(u)c(v)+c(v)c(u)=0,
 \qquad
 c(u)c(v)^*+c(v)^*c(u)=\langle u,v\rangle\Id.
\]
For unit $u$, the operator $N_u=c(u)^*c(u)$ is the projection onto
the states in which $u$ is occupied. If $(e_i)_{i=1}^n$ is an
orthonormal basis, then on $\mathcal H_{p,n}$,
\begin{equation}
 \sum_{i=1}^n N_{e_i}=p\Id.
 \label{ext:number}
\end{equation}
Both assertions can be read directly in the occupation basis.

For a density $\rho$ on $\mathcal H_{p,n}$, define its one-particle
density $\gamma_\rho$ by
\begin{equation}
 a_\rho(u):=\Tr\bigl(c(u)\rho c(u)^*\bigr)
 =\Tr(\rho N_u)=\langle u,\gamma_\rho u\rangle.
 \label{ext:one-body}
\end{equation}
Polarization gives a unique Hermitian operator on $E$. The projection
property of $N_u$ and \eqref{ext:number} imply
\begin{equation}
 0\le\gamma_\rho\le\Id_E,
 \qquad \Tr\gamma_\rho=p.
 \label{ext:pauli-bounds}
\end{equation}
For a Slater input, $\gamma_{\rho_W}=\pi_W$.

\begin{lemma}[One-particle mixture from the original density]
\label{ext:occupation-mixture}
Choose an orthonormal eigenbasis $(f_i)_{i=1}^n$ of $\gamma_\rho$.
For every $p$-element subset $S\subset\{1,\ldots,n\}$, let
$\Omega_S$ be its increasing occupation wedge,
$W_S=\operatorname{span}\{f_i:i\in S\}$, and
$t_S=\langle\Omega_S,\rho\Omega_S\rangle$. Then
\begin{equation}
 t_S\ge0,\qquad \sum_{|S|=p}t_S=1,
 \qquad
 \gamma_\rho=\sum_{|S|=p}t_S\pi_{W_S}.
 \label{ext:diagonal-mixture}
\end{equation}
If $n\ge1$, then for a uniformly distributed unit vector $u\in E$,
any continuous convex $H:[0,1]\to\R$, and any rank-$p$ projection $\pi_W$,
\begin{equation}
 \E_u H(\langle u,\gamma_\rho u\rangle)
 \le \E_u H(\langle u,\pi_Wu\rangle).
 \label{ext:occupation-order}
\end{equation}
If $H$ is strictly convex, equality forces exactly one $t_S$ to be
positive, and consequently $\rho=|\Omega_S\rangle\langle\Omega_S|$.
\end{lemma}
\begin{proof}
The $t_S$ are diagonal entries of the same positive trace-one operator
in a complete orthonormal basis, so they are probabilities. For each
$i$, the occupation operator $N_{f_i}$ is diagonal in this basis and
has eigenvalue $1$ precisely when $i\in S$. Hence
\[
 \langle f_i,\gamma_\rho f_i\rangle
 =\Tr(\rho N_{f_i})=\sum_{S\ni i}t_S.
\]
The off-diagonal entries of $\gamma_\rho$ vanish in the chosen
eigenbasis. This proves \eqref{ext:diagonal-mixture}.
In particular, for every unit $u$ at once,
\[
 a_\rho(u)=\sum_{|S|=p}t_S\langle u,\pi_{W_S}u\rangle.
\]
The projections on the right have the same rank, so their quadratic
readouts have the same spherical distribution. Lemma~\ref{ext:common-mixture}
gives \eqref{ext:occupation-order}.

For $S\ne T$, the Hermitian operator $\pi_{W_S}-\pi_{W_T}$ is
nonzero. Its quadratic form is nonzero somewhere on the unit sphere,
and therefore on a nonempty open set of positive spherical measure.
The strict equality criterion excludes two such positive-weight
indices. The remaining diagonal entry of $\rho$ is one. Positivity
gives
\[
 |\langle\Omega_S,\rho\Omega_T\rangle|^2
 \le t_S t_T
\]
for all occupied subsets, so every other matrix entry is zero.
Thus the original density is the stated pure Slater projector.
\end{proof}

\begin{remark}[The Pauli convex set]
\label{ext:pauli-convex-set}
The set of Hermitian matrices satisfying \eqref{ext:pauli-bounds} is
the convex hull of rank-$p$ orthogonal projections. To see this,
diagonalize such a matrix. The vertices of
$\{x\in[0,1]^n:\sum_i x_i=p\}$ are the zero-one vectors with $p$
ones: a nonintegral point has at least two nonintegral coordinates,
which admit opposite small perturbations preserving the sum.
Rotating a vertex decomposition back proves the assertion.
For the density in the theorem, \eqref{ext:diagonal-mixture} supplies
this decomposition directly from its occupation probabilities and
also yields the full-density equality criterion.
\end{remark}

\subsection{The conditional measurement}

For a unit vector $u$, put
$A_\rho(u)=c(u)\rho c(u)^*$. It is positive with trace
$a_\rho(u)$. The orthogonal splitting
\[
 \Lambda^{p-1}E
 =\Lambda^{p-1}(u^\perp)
 \mathbin{\oplus}
 u\wedge\Lambda^{p-2}(u^\perp)
\]
shows that $\ker c(u)$ in degree $p-1$ is
$\Lambda^{p-1}(u^\perp)$; the second summand is absent when $p=1$.
Since $c(u)^2=0$, the range of $A_\rho(u)$ lies in this kernel.
If $a_\rho(u)>0$, then
\begin{equation}
 \rho_u=\frac{A_\rho(u)}{a_\rho(u)}
 \label{ext:conditional-density}
\end{equation}
is a density on $\Lambda^{p-1}(u^\perp)$.
For a unit Slater vector $\xi$ in that space, adjointness gives
\begin{equation}
 \langle u\wedge\xi,\rho(u\wedge\xi)\rangle
 =\langle\xi,A_\rho(u)\xi\rangle
 =a_\rho(u)\langle\xi,\rho_u\xi\rangle.
 \label{ext:conditioning}
\end{equation}
If $a_\rho(u)=0$, positivity and zero trace give $A_\rho(u)=0$,
so the first expression in \eqref{ext:conditioning} is zero as well.

The relevant Haar decomposition can be written without choosing a
global frame for the bundle $u^\perp$. For $p\ge1$, take
$\Omega_0=e_1\wedge\cdots\wedge e_p$ and embed $U(n-1)$ into
$U(n)$ by $\iota(V)=1\oplus V$. Right invariance gives
\begin{equation}
 \int_{U(n)}\Phi(Q_\rho(U))\,dU
 =\int_{U(n)}\int_{U(n-1)}
 \Phi\bigl(Q_\rho(U\iota(V))\bigr)\,dV\,dU.
 \label{ext:double-haar}
\end{equation}
For fixed $U$, the first vector is $u=Ue_1$, and the remaining wedge
is a Haar Slater vector in $\Lambda^{p-1}(u^\perp)$. The outer
variable $u$ is uniform on the complex unit sphere. Thus
\eqref{ext:conditioning} applies to the same measurement appearing
in \eqref{ext:double-haar}. In particular, the outer measure remains
Haar probability measure; $a_\rho(u)$ enters the argument of $\Phi$.

\begin{lemma}[Slater closure under contraction]
\label{ext:coherent-closure}
If $a_{\rho_W}(u)>0$, the conditional density $(\rho_W)_u$ is a
pure Slater density with occupied space $W\cap u^\perp$.
\end{lemma}
\begin{proof}
Here $a_{\rho_W}(u)=\|\pi_Wu\|^2$. Choose an orthonormal basis of
$W$ whose first vector is $w_1=\pi_Wu/\|\pi_Wu\|$.
Its other vectors $w_2,\ldots,w_p$ lie in $u^\perp$.
Up to the phase of $\Omega_W$, formula~\eqref{ext:contraction} gives
\[
 c(u)\Omega_W
 =\|\pi_Wu\|\,w_2\wedge\cdots\wedge w_p.
\]
Normalizing its rank-one projector proves the assertion, including
the vacuum conditional state when $p=1$.
\end{proof}

\subsection{Induction and equality}

\begin{proof}[Proof of Theorem~\ref{ext:convex-order}]
We induct on $p$, simultaneously for all $n\ge p$. The case $p=0$
has $Q\equiv1$. Suppose $p\ge1$ and the assertion is known in
degree $p-1$. Let $R$ have the Husimi distribution of a Slater input
in $\mathcal H_{p-1,n-1}$ and set $H_\Phi(t)=\E\Phi(tR)$.
This function is continuous and convex. If $\Phi$ is strictly
convex, then $H_\Phi$ is strictly convex, since
\[
 \E R=\binom{n-1}{p-1}^{-1}>0
\]
by \eqref{ext:resolution}.

For every positive-weight conditional density \eqref{ext:conditional-density},
apply the induction hypothesis with the test function
$x\mapsto\Phi(a_\rho(u)x)$. Zero-weight branches give $\Phi(0)$
on both sides. Integrating \eqref{ext:double-haar}, and then applying
Lemma~\ref{ext:occupation-mixture}, yields
\begin{align}
 \int\Phi(Q_\rho)\,d\nu_{p,n}
 &\le\E_u H_\Phi(\langle u,\gamma_\rho u\rangle)\notag\\
 &\le\E_u H_\Phi(\langle u,\pi_Wu\rangle)
 =\int\Phi(Q_{\rho_W})\,d\nu_{p,n}.
 \label{ext:chain}
\end{align}
The final equality follows from Lemma~\ref{ext:coherent-closure}:
each nonzero conditional Slater input has the same lower-dimensional
Husimi distribution. The original double-Haar integrands are continuous
on compact groups. The normalized conditional matrices are measurable
where their continuous weights are positive, and the zero branch was
handled separately. Boundedness therefore justifies these integrations
and their order changes.

If $\Phi$ is strictly convex and the first and last quantities in
\eqref{ext:chain} are equal, the middle Jensen comparison is an
equality. Lemma~\ref{ext:occupation-mixture}, applied to the strictly
convex $H_\Phi$, forces $\rho$ to be a pure Slater density.
Conversely, all Slater inputs have the same Husimi law by Haar
invariance and hence attain equality. The filled sector $p=n$ is
one-dimensional and agrees with this induction.
\end{proof}

\begin{example}[One particle]
\label{ext:one-particle-example}
Every unit vector in $E$ is a one-particle Slater vector. In two
modes, a density with eigenvalues $t$ and $1-t$ has measurement
\[
 Q_\rho(u)=t|u_1|^2+(1-t)|u_2|^2.
\]
At $t=1/2$ this is the constant $1/2$; at $t=1$ it has the
distribution of $|u_1|^2$. Both have mean $1/2$. Theorem~\ref{ext:convex-order}
compares all continuous convex observations between these cases and
every intermediate density. Its strict equality criterion singles
out the rank-one inputs.
\end{example}

\begin{example}[Two-particle coherence]
\label{ext:two-particle-example}
In $\Lambda^2\C^4$, write $e_{ij}=e_i\wedge e_j$ and consider
\[
 \psi=\frac{e_{12}+e_{34}}{\sqrt2},\qquad
 \rho_{\mathrm{mix}}
 =\tfrac12|e_{12}\rangle\langle e_{12}|
  +\tfrac12|e_{34}\rangle\langle e_{34}|.
\]
Both $|\psi\rangle\langle\psi|$ and $\rho_{\mathrm{mix}}$ have
one-particle density $\tfrac12\Id_E$. In the standard eigenbasis,
the shared occupation mixture places weight $1/2$ on each of
$\{1,2\}$ and $\{3,4\}$. Their full Husimi functions differ:
at the plane spanned by $(e_1+e_3)/\sqrt2$ and
$(e_2+e_4)/\sqrt2$, their values are respectively $1/2$ and $1/4$.
Thus the one-particle mixture retains exactly the information needed
for the common conditional-weight comparison, while the conditional
density in \eqref{ext:conditional-density} still carries the original
coherence. Both inputs give strict inequality in
\eqref{ext:order} for every strictly convex test.
\end{example}

\section{Husimi convex order for spin representations}
\label{sec:spin}
\label{spin:section}

An occupation measurement removes one fermionic mode.  Its probability is
an affine function of the Majorana covariance matrix, while its conditional
state lives in a smaller spin representation.  Two facts make this reduction
compatible with convex order: one probability mixture represents the
occupation probabilities in every rotated mode, and a nonzero occupation
projection of a pure Gaussian state remains Gaussian.  We prove both facts
for the original Fock-space action before applying the conditional comparison.

\subsection{Fock space, Majorana operators, and coherent rays}

Let
\[
 \mathcal F_n=\Lambda^\bullet\C^n,
 \qquad
 \mathcal F_n^\epsilon
 =\bigoplus_{k\equiv\epsilon\; (\mathrm{mod}\,2)}\Lambda^k\C^n,
 \qquad \epsilon\in\{0,1\}.
\]
The increasing occupation wedges form an orthonormal basis, denoted by
$|s\rangle$, $s\in\{0,1\}^n$.  Creation is exterior multiplication,
$a_j^*\xi=e_j\wedge\xi$, and annihilation $a_j$ is its Hilbert adjoint.
Thus $a_ja_k+a_ka_j=0$ and
$a_ja_k^*+a_k^*a_j=\delta_{jk}\Id$.  Define
\begin{equation}
 c_{2j-1}=a_j+a_j^*,\qquad c_{2j}=i(a_j^*-a_j).
 \label{spin:majorana}
\end{equation}
These Majorana operators are self-adjoint and satisfy
\begin{equation}
 c_\alpha c_\beta+c_\beta c_\alpha
 =2\delta_{\alpha\beta}\Id.
 \label{spin:clifford}
\end{equation}
For $z\in\C^{2n}$ write $c(z)=\sum_\alpha z_\alpha c_\alpha$.
The quadratic form relevant to \eqref{spin:clifford} is the complex
\emph{bilinear} form
$b(z,w)=\sum_\alpha z_\alpha w_\alpha$.

We use the positive real Clifford realization of $\Spin(2n)$: it consists
of products of an even number of real unit Clifford vectors.  Its Fock
representation sends such a product to the corresponding product of the
$c(v)$, and is denoted by $g\mapsto U_g$.  Each factor $c(v)$ is unitary
and reverses parity; hence $U_g$ is unitary and preserves
$\mathcal F_n^\epsilon$.  Conjugation gives
\begin{equation}
 U_g c(v)U_g^*=c(R_gv),\qquad R_g\in\SO(2n).
 \label{spin:orthogonal-action}
\end{equation}
The map onto $\SO(2n)$ is surjective.  Indeed,
$c(w)c(v)c(w)=c(2\langle w,v\rangle w-v)$ for a real unit $w$;
products of two such reflections give the plane rotations generating
$\SO(2n)$.  Allowing odd products realizes the other component of
$O(2n)$.  An implementing word has odd length precisely when the
orthogonal transformation reverses orientation.  These elementary
Clifford facts also describe the parity change of its Fock action.

A \emph{pure Gaussian ray} in $\mathcal F_n^\epsilon$ means a ray in the
$\Spin(2n)$ orbit of an occupation vector of parity $\epsilon$.
All occupation vectors of that parity belong to the same orbit: changing
an even number of occupation bits is effected, up to a phase, by the
corresponding product of Majorana operators.  Choose a unit representative
$\Omega_{n,\epsilon}$ of this orbit.  These are the coherent rays of the
basic half-spin representation; see also \cite{hackl2021} for their
geometric description.

Each parity representation is irreducible.  To see this directly, the
occupation projections
\[
 |s\rangle\langle s|
 =\prod_{j:s_j=1}a_j^*a_j\prod_{j:s_j=0}(\Id-a_j^*a_j)
\]
belong to the complex span of even Majorana words.  Multiplying them by
even words which change the occupied bits gives every matrix unit
$|t\rangle\langle s|$ with $|s|\equiv|t|\pmod2$.  Thus the algebra
spanned by the representation acts as the full matrix algebra on each
sector.  In particular, with $D=2^{n-1}$ and normalized Haar measure $dg$,
\begin{equation}
 \int_{\Spin(2n)}|U_g\Omega_{n,\epsilon}\rangle
       \langle U_g\Omega_{n,\epsilon}|\,dg
 =\frac{\Id_{\mathcal F_n^\epsilon}}D.
 \label{spin:resolution}
\end{equation}
The average commutes with the representation and has trace one, which
proves the formula.  For a density operator $\rho$ on this sector let
\[
 Q_\rho(g)=\langle U_g\Omega_{n,\epsilon},
                    \rho U_g\Omega_{n,\epsilon}\rangle.
\]
Consequently $0\le Q_\rho\le1$ and $\int Q_\rho\,dg=D^{-1}$.
All parity indices below are interpreted modulo two.

\begin{theorem}[Half-spin Husimi convex order]
\label{spin:convex-order}
Let $n\ge1$ and $\epsilon\in\{0,1\}$.  For every density operator
$\rho$ on $\mathcal F_n^\epsilon$, every continuous convex function
$\Phi:[0,1]\to\R$, and every unit pure Gaussian vector
$\Omega\in\mathcal F_n^\epsilon$,
\begin{equation}
 \int_{\Spin(2n)}\Phi(Q_\rho(g))\,dg
 \le
 \int_{\Spin(2n)}\Phi(Q_{|\Omega\rangle\langle\Omega|}(g))\,dg.
 \label{spin:order}
\end{equation}
The right-hand side is independent of $\Omega$.  If $\Phi$ is strictly
convex, equality holds if and only if $\rho$ is a pure Gaussian projector.
\end{theorem}

\subsection{One covariance decomposition for all rotated modes}

The first-mode occupation projection and the Majorana covariance are
\begin{equation}
 P=a_1^*a_1=\frac{\Id+i c_1c_2}{2},\qquad
 (\Gamma_\rho)_{\alpha\beta}
 =\frac i2\Tr\bigl(\rho[c_\alpha,c_\beta]\bigr).
 \label{spin:covariance}
\end{equation}
The matrix $\Gamma_\rho$ is real and skew-symmetric, and
$\Gamma_{U_g\rho U_g^*}=R_g\Gamma_\rho R_g^{\mathsf T}$.
Putting $u=R_ge_1$ and $v=R_ge_2$, we have
\begin{equation}
 a_\rho(g):=\Tr(PU_g^*\rho U_g)
 =\frac{1+u^{\mathsf T}\Gamma_\rho v}{2}.
 \label{spin:occupation}
\end{equation}
Thus every first-mode probability is a readout of the same covariance.

\begin{lemma}[Covariance mixture within one parity orbit]
\label{spin:covariance-mixture}
For any density operator $\rho$ on $\mathcal F_n^\epsilon$ there are
unit pure Gaussian vectors $\chi_s$ of parity $\epsilon$ and numbers
$t_s\ge0$, $\sum_st_s=1$, such that
\begin{equation}
 \Gamma_\rho=\sum_st_s\Gamma_{\chi_s},
 \qquad
 a_\rho(g)=\sum_st_sa_{\chi_s}(g)\quad\hbox{for every }g\in\Spin(2n).
 \label{spin:shared-mixture}
\end{equation}
The probabilities $t_s$ are independent of $g$.
\end{lemma}
\begin{proof}
Choose $R\in\SO(2n)$ such that
\[
 R\Gamma_\rho R^{\mathsf T}
 =\operatorname{diag}(\lambda_1J,\ldots,\lambda_nJ),
 \qquad J=\begin{pmatrix}0&1\\-1&0\end{pmatrix}.
\]
The usual real skew-symmetric normal form gives an orthogonal $R$.
If its determinant is negative, reversing one coordinate changes the
sign of one $\lambda_j$ and makes $R$ special orthogonal.  The signed
parameters allow this adjustment within the given parity sector.
Choose a Spin lift and let $\sigma$ be the correspondingly rotated
density operator.  Set $t_s=\langle s,\sigma s\rangle$.  Positivity and
trace one give a probability distribution supported on
$|s|\equiv\epsilon\pmod2$.

For an occupation vector,
\[
 \Gamma_{|s\rangle}
 =\operatorname{diag}((2s_1-1)J,\ldots,(2s_n-1)J).
\]
The diagonal two-by-two blocks of $\Gamma_\sigma$ are determined by the
occupation expectations $\Tr(\sigma a_j^*a_j)$; its other blocks vanish
by construction.  It follows that
$\Gamma_\sigma=\sum_st_s\Gamma_{|s\rangle}$.
Rotating the occupation vectors back gives the first identity in
\eqref{spin:shared-mixture}.  Equation~\eqref{spin:occupation} then gives
the second identity simultaneously for every $g$.
\end{proof}

This lemma decomposes a covariance matrix.  The many-body operator
$\rho$ itself need not be a mixture of Gaussian projectors.  What enters
the convex comparison is the common set of probabilities in
\eqref{spin:shared-mixture}.

\subsection{Pure spinors and conditional Gaussian closure}

For a nonzero $\chi\in\mathcal F_n$ define its Clifford annihilator by
\begin{equation}
 L_\chi=\{z\in\C^{2n}:c(z)\chi=0\}.
 \label{spin:annihilator}
\end{equation}
Equation~\eqref{spin:clifford} implies that $L_\chi$ is isotropic for
$b$, so $\dim L_\chi\le n$.  A spinor with $\dim L_\chi=n$ is called
\emph{pure}.  This is an algebraic condition on the original Fock vector.
For conditioning at the last mode, we use
$\mathcal F_0^0=\C$, $\mathcal F_0^1=0$ and call the unique nonzero
zero-mode ray Gaussian.

\begin{lemma}[Gaussian orbit and occupation conditioning]
\label{spin:gaussian-closure}
The pure-spinor rays in $\mathcal F_n^\epsilon$ are exactly its pure
Gaussian rays.  If $b(z,z)=0$, $\chi$ is pure, and $c(z)\chi\ne0$,
then $c(z)\chi$ is pure.  In particular, if $P\chi\ne0$, then under
the occupation-space identification
$P\mathcal F_n^\epsilon=|1\rangle\otimes\mathcal F_{n-1}^{\epsilon-1}$,
\begin{equation}
 \frac{P\chi}{\|P\chi\|}=|1\rangle\otimes\phi
 \label{spin:occupied-closure}
\end{equation}
for a unit pure Gaussian vector $\phi$ in the remaining $n-1$ modes.
For the vacant projection $\Id-P$, the same conclusion holds with
first factor $|0\rangle$ and remaining parity $\epsilon$.
\end{lemma}
\begin{proof}
First let $L\subset\C^{2n}$ be maximal isotropic.  It has trivial
intersection with $\overline L$: a nonzero $z$ in that intersection would
satisfy $z,\overline z\in L$ and
$b(z,\overline z)=\sum_\alpha|z_\alpha|^2>0$.
Choose a Hermitian orthonormal basis
$z_j=(p_j+iq_j)/\sqrt2$ of $L$, with $p_j,q_j$ real.
The equations $b(z_j,z_k)=0$ and
$\sum_\alpha\overline{(z_j)_\alpha}(z_k)_\alpha=\delta_{jk}$ show
that $p_1,q_1,\ldots,p_n,q_n$ form a real orthonormal frame.
An orthogonal transformation carries the standard frame to this one.
Implement it by a product $W$ of real unit Majorana operators, as above.
It carries
$\operatorname{Span}\{e_{2j-1}+ie_{2j}:1\le j\le n\}$ to $L$.
The corresponding standard Clifford operators are $2a_j$.

Their joint kernel is precisely the vacuum line: in an occupation
expansion, $a_j\xi=0$ removes every coefficient of a wedge containing
$e_j$, and requiring this for all $j$ leaves only the scalar wedge.
Conjugating by $W$ proves that the joint kernel associated with $L$ is
the single line $\C W|0\rangle$.
Its parity is the length of the implementing word modulo two.  Thus
two such lines of the same parity are related by an even word, hence
by $\Spin(2n)$.  Conversely, the vacuum has an $n$-dimensional
annihilator, and orthogonal conjugation preserves that dimension.
This proves the equality of the pure-spinor and Gaussian ray sets in
each sector, including the odd reference orbit.

Now suppose $c(z)\chi\ne0$ and $b(z,z)=0$.  Then
$z\notin L_\chi=L_\chi^\perp$, so
$K=L_\chi\cap z^\perp$ has dimension $n-1$.
For $w\in K$, the Clifford relation gives
$c(w)c(z)\chi=0$, and $c(z)^2=0$ gives the same conclusion for $w=z$.
Therefore the annihilator of $c(z)\chi$ contains
$K\oplus\C z$, an $n$-dimensional isotropic space.  The resulting
spinor is pure.

Both $a_1$ and $a_1^*$ are nonzero scalar multiples of isotropic
Clifford vectors.  If $P\chi=a_1^*a_1\chi\ne0$, both successive
actions are nonzero, so $P\chi$ is pure.  Write its normalization as
in \eqref{spin:occupied-closure}.  Its annihilator contains the
$a_1^*$ direction.  Isotropy forces every annihilating vector to have
zero $a_1$ coefficient.  After subtracting its $a_1^*$ component,
the remaining vector lies in the last $2n-2$ Majorana directions.
Modulo the $a_1^*$ line, the annihilator therefore gives an
$(n-1)$-dimensional isotropic annihilator of $\phi$.
The first part of the proof identifies $\phi$ as Gaussian.  Removing
one occupied mode changes its parity from $\epsilon$ to $\epsilon-1$.
For a nonzero vacant projection, use $\Id-P=a_1a_1^*$ and reverse
the two null-vector actions.  Its annihilator contains the $a_1$
direction; quotienting by that line proves the same statement for
the remaining modes, now without a parity change.
\end{proof}

\subsection{Haar conditioning and strict equality}

\begin{proof}[Proof of Theorem~\ref{spin:convex-order}]
We induct on $n$, simultaneously for the two parity sectors.  For $n=1$
each sector is one-dimensional and the assertion is immediate.  For
$n\ge2$ we may choose the reference ray in the form
\[
 \Omega_{n,\epsilon}=|1\rangle\otimes\Omega_{n-1,\epsilon-1}.
\]
Let $G_n=\Spin(2n)$ and let $H_n=\Spin(2n-2)$ be the subgroup
generated by the last $2n-2$ Majorana directions.  Its action on the
occupied first-mode subspace is $\Id\otimes U_h$.  Indeed, every
remaining odd Majorana acquires a first-mode parity sign in the tensor
identification, and the signs cancel in each even word defining $H_n$.
Right invariance of Haar measure gives
\begin{equation}
 \int_{G_n}\Phi(Q_\rho(g))\,dg
 =\int_{G_n}\int_{H_n}
  \Phi\!\left(\langle U_gU_h\Omega_{n,\epsilon},
                  \rho U_gU_h\Omega_{n,\epsilon}\rangle\right)\,dh\,dg.
 \label{spin:double-haar}
\end{equation}
These are integrals of continuous bounded functions on compact groups.

For fixed $g$, identify the positive compression
$P U_g^*\rho U_gP$ with $|1\rangle\langle1|\otimes A_g$.
Its trace is $a_\rho(g)$.  If this trace is positive, define
$\rho_g=A_g/a_\rho(g)$ on $\mathcal F_{n-1}^{\epsilon-1}$.
The measurement value in the inner integral is exactly
\begin{equation}
 a_\rho(g)\,
 \langle U_h\Omega_{n-1,\epsilon-1},
         \rho_gU_h\Omega_{n-1,\epsilon-1}\rangle.
 \label{spin:conditioning}
\end{equation}
If $a_\rho(g)=0$, positivity and trace zero give $A_g=0$, and the
inner integral is $\Phi(0)$.  In both cases the outer measure remains
the original Haar probability $dg$.

Let $R$ have the Husimi distribution of a pure Gaussian state in the
smaller sector, and put $H_\Phi(a)=\E\Phi(aR)$.
This function is continuous and convex on $[0,1]$.
The induction hypothesis, applied for each fixed $a=a_\rho(g)$,
and Lemma~\ref{spin:covariance-mixture} give
\begin{align}
 \int_{G_n}\Phi(Q_\rho)\,dg
 &\le\int_{G_n}H_\Phi(a_\rho(g))\,dg \notag\\
 &\le\sum_st_s\int_{G_n}H_\Phi(a_{\chi_s}(g))\,dg \notag\\
 &=\int_{G_n}H_\Phi(a_{\Omega_{n,\epsilon}}(g))\,dg
 =\int_{G_n}\Phi(Q_{|\Omega_{n,\epsilon}\rangle
                         \langle\Omega_{n,\epsilon}|})\,dg.
 \label{spin:chain}
\end{align}
The middle inequality is the common-mixture Jensen step of
Lemma~\ref{ext:common-mixture}, with probabilities fixed for all $g$.
The next equality follows from Haar invariance of the Gaussian orbit.
The last follows from Lemma~\ref{spin:gaussian-closure}: every nonzero
conditional state of the Gaussian input is a lower-dimensional
Gaussian state, so the inner integral has the same law at each $g$.
Zero branches give $\Phi(0)$ on both sides.

Suppose now that $\Phi$ is strictly convex.  By
\eqref{spin:resolution}, $\E R=2^{-(n-2)}>0$, and hence $H_\Phi$
is strictly convex.  Consider two positive-weight endpoints of the
covariance mixture with difference $\Delta\ne0$.  Choose a real unit
$u$ with $\Delta u\ne0$ and set
$v=-\Delta u/\|\Delta u\|$.  Skew-symmetry implies that $u,v$ are
orthonormal and $u^{\mathsf T}\Delta v\ne0$.
For $n\ge2$, $\SO(2n)$ is transitive on ordered orthonormal
two-frames: complete either frame to an oriented orthonormal basis
and map one basis to the other.  Thus some $g$ distinguishes the two
endpoint probabilities in \eqref{spin:occupation}.  Continuity
preserves this difference on a nonempty open set, which has positive
Haar measure.  Jensen's inequality in \eqref{spin:chain} is then
strict.

Equality therefore forces all positive-weight endpoints in the
construction of Lemma~\ref{spin:covariance-mixture} to have the same
covariance.  Distinct occupation strings have distinct diagonal
covariances, so exactly one $t_s$ is nonzero.  For a positive operator,
\[
 |\langle s,\sigma t\rangle|^2
 \le\langle s,\sigma s\rangle\langle t,\sigma t\rangle;
\]
this is Cauchy--Schwarz applied to $\sigma^{1/2}|s\rangle$ and
$\sigma^{1/2}|t\rangle$.  The unique nonzero diagonal entry, together
with trace one, forces $\sigma=|s\rangle\langle s|$.
Undoing the Spin rotation shows that $\rho$ is a pure Gaussian
projector.  Such projectors attain equality by Haar invariance.
\end{proof}

\begin{example}[Occupation conditioning and shared probabilities]
\label{spin:example}
In two modes,
$\chi_\theta=\cos\theta\,|00\rangle+\sin\theta\,|11\rangle$
is Gaussian: it is obtained from the vacuum by
$\exp[\theta(a_1^*a_2^*-a_2a_1)]$, whose generator equals
$(c_1c_3-c_2c_4)/2$.
Its first-mode occupation probability is $\sin^2\theta$;
whenever this is nonzero, the remaining conditional state is the
one-mode occupied ray.

In four modes, consider
$\psi=(|0000\rangle+|1111\rangle)/\sqrt2$.
Its covariance is zero: each mode is occupied with probability
$1/2$, and a quadratic Majorana expression cannot connect its two
occupation components.  Thus its covariance mixture uses the two
fixed weights $1/2,1/2$ on the vacuum and the fully occupied state.
For every rotated mode, \eqref{spin:shared-mixture} gives
$a_\psi(g)=1/2$.  The state is non-Gaussian, since a Gaussian
covariance is an orthogonal conjugate of the vacuum covariance
and therefore squares to $-\Id$.  Its first-mode occupied
conditional state is nevertheless the Gaussian vector $|111\rangle$.
This example separates covariance mixing, conditional states, and
the original many-body state.
\end{example}

\subsection{The basic odd-spin representation}

The odd-dimensional spin representation has the same coherent rays
as a half-spin representation in one larger real dimension.  Equality
of the invariant probability measures on those rays transfers the
distribution comparison without altering its normalization.

\begin{corollary}[Odd-spin Husimi convex order]
\label{spin:odd}
For every $m\ge1$, the basic $2^m$-dimensional representation of
$\Spin(2m+1)$ satisfies the conclusion of
Theorem~\ref{spin:convex-order} on its highest-weight coherent orbit.
For a strictly convex test function, equality holds exactly at the
pure coherent projectors of that orbit.
\end{corollary}
\begin{proof}
Set $N=m+1$.  Fix a real unit vector $e\in\R^{2N}$ and embed
$H=\Spin(2N-1)$ as the lift of its stabilizer in $\SO(2N)$.
The restriction of either half-spin representation to $H$ is the
basic spin representation.  One algebraic verification uses
\[
 \operatorname{Cl}_{2N-1}^0(\C)
 \simeq\operatorname{Cl}_{2N-2}(\C)
 \simeq\operatorname{Mat}_{2^{N-1}}(\C).
\]
The induced unital action on a half-spin space of dimension
$2^{N-1}$ is the standard irreducible matrix-module action.
The Clifford identities follow by using products with one fixed
unit generator for the first isomorphism and the creation and
annihilation matrix units for the second.  Under a maximal torus
of $H$, the weights are
$(\pm\tfrac12,\ldots,\pm\tfrac12)$, each with multiplicity one:
they are the restrictions of the half-spin weights, with the last
sign determined by parity.  In particular, a highest-weight ray
is an occupation Gaussian ray.

For completeness we identify the entire orbit.  The proof of
Lemma~\ref{spin:gaussian-closure} identifies maximal isotropic
annihilators with orthogonal complex structures on $\R^{2N}$.
For example, with $z_j=(p_j+iq_j)/\sqrt2$ one may take
$Jp_j=q_j$ and $Jq_j=-p_j$.  This construction is independent
of the unitary basis of the isotropic space, and the two parity
orbits correspond to its two possible orientation components.

Let $J_1,J_2$ lie in the same component.  An element of
$\SO(2N-1)$ fixing $e$ can first send $J_1e$ to $J_2e$.
After this conjugation the two complex structures agree on
$\operatorname{Span}(e,J_2e)$.  On its orthogonal complement
they induce orthogonal complex structures of the same orientation.
Choose adapted oriented orthonormal bases there; the map between
these bases belongs to $\SO(2N-2)$ and conjugates one restricted
complex structure to the other.  Extending it by the identity on
$\operatorname{Span}(e,J_2e)$ fixes $e$ and completes the conjugation.
Thus $H$ is transitive on the entire half-spin coherent orbit.

The pushforwards of Haar probability from $\Spin(2N)$ and from
$H$ coincide: both are invariant probabilities on this compact
transitive $H$-space.  Uniqueness follows by averaging any
continuous function over $H$; the average is constant on the
orbit, so every invariant probability has the same integral.
The Husimi functions are those of the same rays in the same
Hilbert space.  Theorem~\ref{spin:convex-order}, with $n=N$,
therefore gives the asserted inequality and equality criterion.
\end{proof}

\section{Coherent distributions and entropy minima}
\label{ent:section}

\subsection{Normalization and entropy conventions}
\label{ent:definitions}

In this section, $(\mathcal O,\nu)$ denotes either the Slater orbit
with its invariant probability measure from Section~\ref{ext:section},
or a fixed-parity Gaussian orbit from Theorem~\ref{spin:convex-order}.
The representation dimension is respectively
$D=\binom np$ or $D=2^{n-1}$. The same definitions apply to the odd-spin
orbit in Corollary~\ref{spin:odd}. A unit coherent vector at $x\in\mathcal O$
is denoted by $\Omega_x$, and
$Q_\rho(x)=\langle\Omega_x,\rho\Omega_x\rangle$ is independent of its
phase.

The resolution identity gives
\begin{equation}
 \int_{\mathcal O}|\Omega_x\rangle\langle\Omega_x|\,d\nu(x)
 =\frac{\Id}{D},
 \qquad D\int_{\mathcal O}Q_\rho\,d\nu=1.
 \label{ent:resolution}
\end{equation}
For exterior powers this is \eqref{ext:resolution}. For the irreducible
spin representations, the averaged projector commutes with the group;
Schur's lemma and its trace one give the identity.

Define the unnormalized and normalized positive moments by
\begin{equation}
 M_\rho(q)=\int_{\mathcal O}Q_\rho^q\,d\nu,
 \qquad N_\rho(q)=D M_\rho(q),\qquad q>0.
 \label{ent:moments-definition}
\end{equation}
The normalized first moment is one. Since $Q_\rho\ge0$ has positive
integral, each $M_\rho(q)$ is strictly positive. All these moments are
finite because $0\le Q_\rho\le1$.

We use the resolution measure $D\,d\nu$ to define the
R\'enyi--Wehrl and Wehrl entropies:
\begin{align}
 S_q(\rho)&=\frac{\log N_\rho(q)}{1-q},
 &&q>0,\quad q\ne1,
 \label{ent:renyi-definition}\\
 S_W(\rho)&=-D\int_{\mathcal O}Q_\rho\log Q_\rho\,d\nu.
 \label{ent:wehrl-definition}
\end{align}
Logarithms are natural and $0\log0=0$. In particular, the latter
integrand is the continuous extension of $x\log x$ at zero.
The probability density relative to $\nu$ is $f_\rho=DQ_\rho$;
its differential entropy is
\begin{equation}
 -\int f_\rho\log f_\rho\,d\nu=S_W(\rho)-\log D.
 \label{ent:output-entropy}
\end{equation}

\subsection{The coherent Beta laws}

For $a,b>0$, the distribution $\operatorname{Beta}(a,b)$ on $(0,1)$
has density
\[
 \frac{\Gamma(a+b)}{\Gamma(a)\Gamma(b)}
 t^{a-1}(1-t)^{b-1},
\]
where $\Gamma(z)=\int_0^\infty s^{z-1}e^{-s}\,ds$ for $z>0$
is Euler's Gamma function. A $\operatorname{Gamma}(a,1)$ variable
has density $s^{a-1}e^{-s}/\Gamma(a)$ on $(0,\infty)$.
If $X$ and $Y$ are independent Gamma variables with shapes $a,b$
and common scale one, the change of variables
$(X,Y)=(rt,r(1-t))$ shows that $X/(X+Y)$ has distribution
$\operatorname{Beta}(a,b)$.

\begin{proposition}[Coherent measurement laws]
\label{ent:beta-products}
For a Slater input in $\mathcal H_{p,n}$ with $0<p<n$,
\begin{equation}
 Q_{\rho_W}\ \overset{\mathrm{law}}=
 \prod_{j=0}^{p-1} A_j,
 \qquad
 A_j\sim\operatorname{Beta}(p-j,n-p),
 \label{ent:exterior-beta}
\end{equation}
where the $A_j$ are independent. For $p=0$ or $p=n$, the coherent
measurement is identically one.

For a pure Gaussian input in either parity sector $\mathcal F_n^\epsilon$,
\begin{equation}
 Q_{\Omega}\ \overset{\mathrm{law}}=
 \prod_{k=1}^{n-1} B_k,
 \qquad B_k\sim\operatorname{Beta}(k,k),
 \label{ent:spin-beta}
\end{equation}
with independent factors. The product is one when $n=1$.
For the basic spin representation of $\Spin(2m+1)$, use
\eqref{ent:spin-beta} with $n=m+1$.
\end{proposition}
\begin{proof}
Let $Z$ be a standard complex Gaussian vector in $E=\C^n$, with
independent coordinates of density $\pi^{-1}e^{-|z|^2}$. Its direction
$u=Z/\|Z\|$ is uniform on the complex unit sphere. For a fixed
$p$-plane $W$, the variables
\[
 X=\|\pi_W Z\|^2,\qquad
 Y=\|(\Id-\pi_W)Z\|^2
\]
are independent Gamma variables of shapes $p$ and $n-p$.
Consequently the first conditional weight is
\begin{equation}
 a_{\rho_W}(u)=\|\pi_Wu\|^2=\frac{X}{X+Y}
 \sim\operatorname{Beta}(p,n-p).
 \label{ent:first-exterior-weight}
\end{equation}
For almost every $u$ this weight is positive. By
Lemma~\ref{ext:coherent-closure}, the normalized conditional input
is a Slater density in $\Lambda^{p-1}(u^\perp)$. Its inner Haar
Husimi distribution depends only on $(p-1,n-1)$, independently of
$u$. On the double-Haar probability space of
\eqref{ext:double-haar}, this conditional distribution being constant
means that the lower-dimensional readout is independent of the first
weight. Iterating \eqref{ext:conditioning} proves
\eqref{ent:exterior-beta}; at step $j$ the two Beta shapes are
$p-j$ and $(n-j)-(p-j)=n-p$.

For a Gaussian input, its Majorana covariance $\Gamma_\Omega$
satisfies $\Gamma_\Omega^2=-\Id$. A Haar-rotated first mode gives
a uniform oriented orthonormal two-frame $(u,v)$ in $\R^{2n}$.
Conditioned on $u$, the vector $v$ is uniform on the unit sphere in
$u^\perp$. The linear functional $v\mapsto u^{\mathsf T}\Gamma_\Omega v$
is represented there by the unit vector $\Gamma_\Omega^{\mathsf T}u$.
For $n\ge2$, its value $t$ therefore has density proportional to
\[
 (1-t^2)^{n-2},\qquad -1<t<1.
\]
The occupation identity \eqref{spin:occupation} gives
$a=(1+t)/2$, whose density is proportional to
$a^{n-2}(1-a)^{n-2}$. Thus
$a\sim\operatorname{Beta}(n-1,n-1)$.
The Gaussian conditional closure in Lemma~\ref{spin:gaussian-closure}
and the double-Haar identity \eqref{spin:double-haar} make the
remaining coherent readout independent of this weight, with the
law for $n-1$ modes. Induction proves \eqref{ent:spin-beta}.
The one-dimensional initial sector gives the empty product.
Finally, Corollary~\ref{spin:odd} identifies the entire odd-spin
coherent orbit and its probability measure with the indicated
half-spin orbit.
\end{proof}

The independence in Proposition~\ref{ent:beta-products} describes
coherent inputs. For a general density, its normalized conditional
state varies with the preceding measurement direction; the convex-order
proof retains that dependence.

\subsection{Positive moments and explicit minima}

The Beta integral gives
\begin{equation}
 \E A^q=
 \frac{\Gamma(a+q)\Gamma(a+b)}
 {\Gamma(a)\Gamma(a+b+q)},
 \qquad A\sim\operatorname{Beta}(a,b),\quad q>0.
 \label{ent:beta-moment}
\end{equation}
Consequently the coherent moments in the two series are
\begin{align}
 m_{p,n}(q)
 &:=\prod_{j=0}^{p-1}
 \frac{\Gamma(p-j+q)\Gamma(n-j)}
 {\Gamma(p-j)\Gamma(n-j+q)},
 \label{ent:exterior-moments}\\
 b_n(q)
 &:=\prod_{k=1}^{n-1}
 \frac{\Gamma(k+q)\Gamma(2k)}
 {\Gamma(k)\Gamma(2k+q)}.
 \label{ent:spin-moments}
\end{align}
Here $m_{0,n}(q)=1$ by the empty-product convention, while each
factor in $m_{n,n}(q)$ cancels to one. These formulas cover every
real $q>0$. In particular,
\begin{equation}
 m_{p,n}(1)=\binom np^{-1},\qquad
 b_n(1)=2^{-(n-1)},
 \label{ent:first-moments}
\end{equation}
as required by \eqref{ent:resolution}.

\begin{theorem}[Wehrl and R\'enyi--Wehrl minima]
\label{ent:minima}
Let $H_0=0$ and $H_r=\sum_{j=1}^r j^{-1}$ for $r\ge1$.
For every density on $\mathcal H_{p,n}$, $0\le p\le n$,
\begin{align}
 S_q(\rho)
 &\ge\frac{\log\!\left[\binom np\,m_{p,n}(q)\right]}{1-q},
 &&q>0,\quad q\ne1,
 \label{ent:exterior-renyi}\\
 S_W(\rho)
 &\ge s_{p,n}:=\sum_{j=0}^{p-1}(H_{n-j}-H_{p-j}).
 \label{ent:exterior-wehrl}
\end{align}
Equality in each inequality holds precisely for pure Slater densities.

For every density supported on the fixed parity sector
$\mathcal F_n^\epsilon$, $n\ge1$,
\begin{align}
 S_q(\rho)
 &\ge\frac{\log[2^{n-1}b_n(q)]}{1-q},
 &&q>0,\quad q\ne1,
 \label{ent:spin-renyi}\\
 S_W(\rho)
 &\ge s_n^{\mathrm{spin}}:=\sum_{k=1}^{n-1}(H_{2k}-H_k).
 \label{ent:spin-wehrl}
\end{align}
Equality holds precisely for pure Gaussian densities in that parity
sector. For the basic spin representation of $\Spin(2m+1)$, the
same formulas hold with dimension $2^m$, moment $b_{m+1}(q)$,
and constant $s_{m+1}^{\mathrm{spin}}$.
\end{theorem}
\begin{proof}
For $q>1$, the function $x\mapsto x^q$ is continuous and strictly
convex on $[0,1]$. Theorems~\ref{ext:convex-order} and
\ref{spin:convex-order} imply that its moment is maximized by a
pure coherent input. Multiplication by $D$, application of the
increasing logarithm, and division by $1-q<0$ give the entropy
lower bound. For $0<q<1$, use the continuous strictly convex
function $x\mapsto-x^q$. The moment inequality reverses, and
$1-q>0$, giving the same direction for the entropy inequality.
The positivity of the moments justifies every logarithm and its
equality criterion. Strict convexity identifies the equality cases
in both ranges.

For Wehrl entropy, apply the same theorems to the continuous strictly
convex function $x\mapsto x\log x$, with its value zero at the
origin, and multiply by $-D$. It remains to evaluate the coherent
integral. Differentiation of a moment at $q=1$ gives
\[
 \left.\frac{d}{dq}\int Q_\Omega^q\,d\nu\right|_{q=1}
 =\int Q_\Omega\log Q_\Omega\,d\nu.
\]
This differentiation is justified directly: for
$q\in[1/2,3/2]$ and $0<x\le1$,
$|x^q\log x|\le x^{1/2}|\log x|\le2/e$.
The derivative at $x=0$ is zero, so dominated convergence applies
on the probability space including all zeros of $Q_\Omega$.

Taking logarithmic derivatives in
\eqref{ent:exterior-moments} and \eqref{ent:spin-moments}, and
using $\Gamma(z+1)=z\Gamma(z)$, gives
\begin{align*}
 -D_{p,n}m_{p,n}'(1)
 &=\sum_{j=0}^{p-1}(H_{n-j}-H_{p-j}),\\
 -2^{n-1}b_n'(1)
 &=\sum_{k=1}^{n-1}(H_{2k}-H_k).
\end{align*}
For clarity, if $\psi_\Gamma=\Gamma'/\Gamma$, the recurrence
$\psi_\Gamma(r+1)=\psi_\Gamma(r)+1/r$ shows that each difference
$\psi_\Gamma(b+1)-\psi_\Gamma(a+1)$ at positive integers equals
$H_b-H_a$. Equation~\eqref{ent:first-moments} supplies the factor
$D$ in the preceding derivatives. This proves the displayed constants.
The odd-spin conclusions follow from Corollary~\ref{spin:odd}.
\end{proof}

\begin{example}[Small sectors]
\label{ent:examples}
For one particle in $n$ modes, the coherent Wehrl value is
$s_{1,n}=H_n-1$. For two particles in four modes,
\[
 s_{2,4}=(H_4-H_2)+(H_3-H_1)=\frac{17}{12}.
\]
The maximally mixed density on this six-dimensional sector has
$Q\equiv1/6$ and hence $S_W=\log6$, larger than the coherent
value. In the spin series the constants for $n=1,2,3$ are
$0$, $1/2$, and $13/12$. The empty and filled exterior sectors
have zero Wehrl and R\'enyi--Wehrl entropy.
\end{example}

All formulas above concern the basic exterior, half-spin, and odd-spin
representations specified in the corresponding theorems. The
representation parameter is not replaced by a higher Cartan power.

\section{A polynomial obstruction to three Gaussian summands}
\label{sec:gaussian-rank}

A Gaussian decomposition is an exact vector identity. Its coefficients may
have arbitrary complex phases, and its summands may pair modes belonging to
different prescribed blocks. Consequently, estimates on each summand in
isolation do not determine the minimum number of summands. We use a
polynomial of one complete Clifford Gram operator: a common change of
coordinates transports every summand and every Gram entry simultaneously.

Let $\mathcal F_n=\Lambda\C^n$ and
$\mathcal F_n^+=\Lambda^{\mathrm{even}}\C^n$.
Products of vectors in an exterior algebra below mean wedge products.
For an ordered direct sum $E_A\oplus E_B$, we identify
$\Lambda E_A\otimes\Lambda E_B$ with $\Lambda(E_A\oplus E_B)$ by
$s\otimes t\mapsto s\wedge t$, with all $A$ modes before the $B$ modes.
Denote by
$\widehat{\mathcal G}_n^+$ the nonzero complex cone over the vacuum orbit of
the compact Euclidean group $\Spin(2n)$. For $0\ne\psi\in\mathcal F_n^+$ set
\begin{equation}
 \chi_G(\psi)=\min\left\{r:\psi=g_1+\cdots+g_r,
                 \quad g_j\in\widehat{\mathcal G}_n^+\right\}.
 \label{rank:definition}
\end{equation}
Every even occupation vector belongs to this cone, so the minimum exists.
Allowing odd Gaussian summands gives the same minimum for an even target:
the sum of all odd summands is zero, and deleting them preserves the vector
without increasing the number of terms. Nonzero scalar multiplication also
preserves \eqref{rank:definition}, in both directions.

\begin{theorem}\label{rank:main}
For any nonzero $\psi,\phi\in\mathcal F_4^+$,
\begin{equation}
 \chi_G(\psi\otimes\phi)=\chi_G(\psi)\chi_G(\phi).
 \label{rank:multiplicativity}
\end{equation}
Each factor has rank one or two. If both factors are non-Gaussian, their
product has rank four. In particular, for
\[
 M=\frac{|0000\rangle+|1111\rangle}{\sqrt2},
 \qquad \chi_G(M^{\otimes2})=4.
\]
The decomposition on the left permits arbitrary eight-mode Gaussian terms,
including terms with pairing across the two four-mode blocks.
\end{theorem}

The last assertion is the unrestricted two-copy rank question posed by
Cudby--Strelchuk~\cite{cudby2023}. Gaussian fidelity and extent, for which
four-mode multiplicativity results were obtained in
\cite{diaskoenig2024,reardon2024}, are different quantities: they retain
normalization or coefficient costs, whereas \eqref{rank:definition} counts
nonzero terms in an exact sum.

\subsection{Clifford conventions and the physical cone}

Write $E=\C^n$, $U=E^*\oplus E$, and
\begin{equation}
 c(f,v)=\iota_f+\varepsilon_v,\qquad
 \varepsilon_v s=v\wedge s,\qquad q(f,v)=f(v).
 \label{rank:clifford}
\end{equation}
The polar form is
$b((f,v),(g,w))=f(w)+g(v)$, and
\begin{equation}
 c(u)^2=q(u)\Id,\qquad
 c(u)c(w)+c(w)c(u)=b(u,w)\Id.
 \label{rank:car}
\end{equation}
A nonzero $s\in\mathcal F_n$ is pure if
$\Ann(s)=\{u:c(u)s=0\}$ has dimension $n$.
The CAR make this annihilator isotropic, so its dimension is at most $n$.
The common kernel of all contractions is the scalar vacuum line. Indeed,
if $s=\sum_I s_Ie_I$ has a nonempty coefficient $s_I$, contracting by any
$i\in I$ reads $s_I$, up to sign, in the $I\setminus\{i\}$ coefficient;
no other source coefficient contributes to that same output.

By Lemma~\ref{spin:gaussian-closure}, even pure-spinor rays are exactly the
compact Gaussian vacuum rays; see also~\cite{hackl2021}.
To specify the real form in the split coordinates, use the standard
Hermitian basis of $E$ and define
\[
 \overline{(f,v)}=(v^\flat,f^\sharp),\qquad
 v^\flat(e_i)=\overline{v_i},\qquad
 f^\sharp=\sum_i\overline{f(e_i)}e_i.
\]
This exchanges the two split summands and is not coordinatewise
conjugation on $E^*\oplus E$. It satisfies
$b(u,\overline u)>0$ for $u\ne0$. A maximal isotropic subspace $L$ therefore
has $L\cap\overline L=0$. An orthonormal basis of $L$ gives, by taking real
and imaginary parts, a real orthogonal frame, and an orthogonal map sends
the vacuum annihilator to $L$. A lift to the Pin group transports the
corresponding annihilator. Its spinor line is unique, because transport
back to the vacuum reduces uniqueness to the common contraction kernel.
If that line is even, the lift has even Clifford parity and lies in Spin.
Thus the compact orbit gives every even pure ray. The converse follows
from the CAR transport of the annihilator. In particular,
\begin{equation}
 \widehat{\mathcal G}_n^+
 =\{s\in\mathcal F_n^+:s\ne0,\ \dim\Ann(s)=n\}.
 \label{rank:cone}
\end{equation}
Consequently complex Spin transformations preserve the same physical cone
of rays. They are used below as invertible coordinate transformations,
not as norm-preserving physical protocols.

For later use, a nonzero isotropic Clifford multiplication preserves
purity. If $s$ is pure, $q(u)=0$, and $c(u)s\ne0$, then $u\notin\Ann(s)$
and
\[
 \C u+\bigl(\Ann(s)\cap u^\perp\bigr)\subseteq\Ann(c(u)s).
\]
The left side has dimension $n$: the restriction of $b(u,\cdot)$ to the
maximal isotropic space $\Ann(s)$ is nonzero, and $u$ does not belong to
that space. The upper dimension bound therefore gives purity. Two such
multiplications prove the closure of a nonzero occupation or vacancy
projection, consistent with Lemma~\ref{spin:gaussian-closure}.

\subsection{The whole vacuum chart from elementary units}

Fix an ordered basis $e_0,\ldots,e_{n-1}$ and dual basis
$f_0,\ldots,f_{n-1}$. Write $\varepsilon_i=\varepsilon_{e_i}$ and
$\iota_i=\iota_{f_i}$. Let $[s]_0$ denote the vacuum coefficient and
$[s]_2\in\Lambda^2E$ the degree-two homogeneous component.
Suppose $s$ is pure and $a=[s]_0\ne0$. Projection of $\Ann(s)$ to $E^*$
is injective: if $\varepsilon_v s=0$, its degree-one part is $av=0$.
Since both spaces have dimension $n$, the projection is an isomorphism.
It defines a linear map $\mathcal D:E^*\to E$ with
\begin{equation}
 (\iota_f+\varepsilon_{\mathcal D f})s=0\quad(f\in E^*),\qquad
 f(\mathcal D g)+g(\mathcal D f)=0.
 \label{rank:graph}
\end{equation}
The second identity follows from the CAR. Applying a second contraction
and then reading degree zero gives
\begin{equation}
 g(\mathcal D f)=-\frac{[\iota_g\iota_fs]_0}{a}.
 \label{rank:graph-coordinates}
\end{equation}
These are equations on the complete original vector, not just its first
two homogeneous components.

\begin{lemma}[Elementary chart reconstruction]\label{rank:chart}
For $i<j$ define
\[
 a_{ij}=f_i(\mathcal D f_j)=-f_j(\mathcal D f_i),\qquad
 \mathcal U(\mathbf a)=\prod_{i<j}(\Id+a_{ij}\varepsilon_i\varepsilon_j),
 \qquad
 \Omega(\mathbf a)=\prod_{i<j}(1+a_{ij}e_ie_j).
\]
Here $\mathbf a=(a_{ij})_{i<j}$ is the coefficient array, distinct from
the vacuum coefficient $a$. The factors may be taken in any fixed order.
Then $s=a\Omega(\mathbf a)$.
More generally, vectors annihilated by the same graph
$\iota_f+\varepsilon_{\mathcal D f}$ and having the same vacuum coefficient
are equal.
\end{lemma}
\begin{proof}
Put $N_{ij}=\varepsilon_i\varepsilon_j$. The creation CAR give
$N_{ij}^2=0$, $[N_{ij},\varepsilon_v]=0$, and
$[N_{ij},N_{k\ell}]=0$, including repeated-index cases.
Thus $U_{ij}(z)=\Id+zN_{ij}$ has inverse $\Id-zN_{ij}$.
For $h_{ij}(f)=f(e_i)e_j-f(e_j)e_i$, the mixed CAR give
\[
 [\iota_f,N_{ij}]=\varepsilon_{h_{ij}(f)},\qquad
 N_{ij}\varepsilon_{h_{ij}(f)}
 =\varepsilon_{h_{ij}(f)}N_{ij}=0.
\]
It follows, on all of $\mathcal F_n$, that
\begin{equation}
 U_{ij}(z)^{-1}\iota_fU_{ij}(z)
 =\iota_f+z\varepsilon_{h_{ij}(f)}.
 \label{rank:elementary-shear}
\end{equation}
Set $h_{\mathbf a}(f)=\sum_{i<j}a_{ij}h_{ij}(f)$.
Evaluating on the dual basis and using \eqref{rank:graph}, including its
zero diagonal, gives $h_{\mathbf a}=-\mathcal D$.
All factors commute with every creator, so composing
\eqref{rank:elementary-shear} gives
\[
 \iota_f\mathcal U(\mathbf a)^{-1}s
 =\mathcal U(\mathbf a)^{-1}(\iota_f+\varepsilon_{\mathcal D f})s=0.
\]
Each factor and its inverse preserve degree zero. The common contraction
kernel is the vacuum line, hence
$\mathcal U(\mathbf a)^{-1}s=a\,1$, and multiplication by
$\mathcal U(\mathbf a)$ proves
the assertion. Applying the same inverse product to two vectors proves
uniqueness. No expansion of their higher-degree components is needed.
\end{proof}

For reference, if $\omega=\sum_{i<j}a_{ij}e_ie_j$, then
$\Omega(\mathbf a)=\exp(\omega)$: the degree-two monomials commute and square to
zero, so their finite exponentials multiply to the displayed product.
The proof of Lemma~\ref{rank:chart} does not require a general exponential
construction.

For $g\in\GL(E)$ the original exterior map satisfies
\begin{equation}
 \iota_f\Lambda g=\Lambda g\,\iota_{f\circ g},\qquad
 \varepsilon_{gv}\Lambda g=\Lambda g\,\varepsilon_v.
 \label{rank:gl-intertwining}
\end{equation}
The first identity follows by induction on exterior monomials from the
contraction rule, and the second from multiplicativity.
Write $\Omega_{\mathcal D}$ for the chart product with coefficients
$a_{ij}=f_i(\mathcal D f_j)$, and let $g^\vee:E^*\to E^*$ denote the
algebraic dual map $g^\vee f=f\circ g$, with no complex conjugation.
The new graph is
\begin{equation}
 \mathcal D^g(f)=g\bigl(\mathcal D(f\circ g)\bigr).
 \label{rank:graph-gl}
\end{equation}
Both $\Lambda g\Omega_{\mathcal D}$ and $\Omega_{\mathcal D^g}$ have this graph and vacuum
coefficient one. Lemma~\ref{rank:chart} therefore proves their equality.
In particular a nondegenerate alternating graph can be reduced to its
standard form by a single $g$ acting on the whole Fock space.

Here is an explicit linear-algebra justification of that reduction.
For a nondegenerate alternating form choose $u,v$ with pairing one.
Their span is nondegenerate. Subtracting the appropriate multiples of
$u,v$ from any vector makes it orthogonal to both, giving the direct sum
of this plane and its orthogonal complement. The complement is
nondegenerate, since a vector orthogonal to both summands is orthogonal
to the whole space. Repeating produces a symplectic basis. Applied to the
alternating tensor $\mathcal D:E^*\to E$, this is exactly the congruence
$\mathcal D\mapsto g\mathcal D g^\vee$ in \eqref{rank:graph-gl}.

\begin{example}\label{rank:four-mode-example}
The four-mode vector $1+\kappa e_0e_1e_2e_3$, with $\kappa\ne0$, is not
Gaussian. Its degree-two component vanishes, so
\eqref{rank:graph-coordinates} would force $\mathcal D=0$ if it were pure;
Lemma~\ref{rank:chart} would then make it a vacuum multiple. It is a sum
of two Gaussian occupation vectors, and hence has rank two.
Its two-copy expansion has four terms, but this observation alone does
not rule out a shorter expansion using eight-mode Gaussians.
\end{example}

\subsection{A relative invariant of the complete Clifford Gram}

For the remainder of the rank-four proof let $n=8$ and
$\mathrm{vol}=e_0\cdots e_7$. Fix the bilinear Chevalley convention
\begin{equation}
 \mathcal B(s,t)=[s\wedge\rev(t)]_{\mathrm{vol}},\qquad
 \rev|_{\Lambda^jE}=(-1)^{j(j-1)/2}\Id.
 \label{rank:chevalley}
\end{equation}
Reversal acts on the second leg; no complex conjugation is present.
For an ordered subset $I\subseteq\{0,\ldots,7\}$,
\begin{equation}
 \mathcal B(e_I,e_{I^c})=(-1)^{\sum_{i\in I}(i+1)}.
 \label{rank:chevalley-sign}
\end{equation}
Indeed, for $|I|=k$ the wedge-order sign has exponent
$\sum_{i\in I}i-k(k-1)/2$, and reversing the second leg adds
$(8-k)(7-k)/2$. Their sum is congruent to
$\sum_{i\in I}(i+1)$ modulo two. This calculation applies in every degree.
The form is symmetric and nondegenerate.

For every Clifford direction $u$, the exterior multiplication and
contraction rules give
\begin{equation}
 \mathcal B(c(u)s,t)=-\mathcal B(s,c(u)t).
 \label{rank:bilinear-adjoint}
\end{equation}
For creation, move the created vector past the other seven factors in a
nonzero top-degree term; for contraction, apply contraction to the
identically zero degree-nine exterior product. Both yield the displayed
minus sign with \eqref{rank:chevalley}. Consequently an elementary
creation unit, or its two-contraction counterpart, preserves $\mathcal B$:
its generator is skew-adjoint for $\mathcal B$ and squares to zero.
Also $\tau_i=\varepsilon_i+\iota_{f_i}$ has $\tau_i^2=\Id$ and multiplier
$-1$, so an even product of the $\tau_i$ preserves $\mathcal B$.
Finally,
\begin{equation}
 \mathcal B(\Lambda g\,s,\Lambda g\,t)
       =\det(g)\mathcal B(s,t).
 \label{rank:gl-multiplier}
\end{equation}
This follows by applying $\Lambda g$ to the top exterior product.

Let $W=\Lambda^2U$ and equip it with
\[
 G(u\wedge v,w\wedge z)=b(u,w)b(v,z)-b(u,z)b(v,w).
\]
Define the whole bivector source map and its raised Gram by
\begin{equation}
 J_\psi(u\wedge v)=\tfrac12[c(u),c(v)]\psi,\qquad
 H_\psi(\xi,\zeta)=\mathcal B(J_\psi\xi,J_\psi\zeta),\qquad
 A_\psi=G^{-1}H_\psi.
 \label{rank:full-gram}
\end{equation}
All $\binom{16}{2}=120$ bivector directions occur here.
The matrix of $A_\psi$ is $G^{-1}J_\psi^{\mathsf T}\mathcal BJ_\psi$;
it need not be Hermitian or positive. Set
\begin{align}
 s_\psi&=\tfrac12\mathcal B(\psi,\psi),&
 t_k(\psi)&=\Tr(A_\psi^k),\label{rank:moments}\\
 F_{16}(\psi)&=576s_\psi^2t_6(\psi)-7t_4(\psi)^2
 -18960s_\psi^4t_4(\psi)+4205376s_\psi^8,&
 P_{18}(\psi)&=s_\psi F_{16}(\psi).
 \label{rank:polynomial}
\end{align}
These are homogeneous polynomials of degrees $16$ and $18$ in $\psi$.

\begin{lemma}\label{rank:covariance}
Suppose invertible maps $T$ on $\mathcal F_8$ and $R$ on $U$ satisfy
\[
 Tc(u)T^{-1}=c(Ru),\qquad R^*b=b,\qquad
 \mathcal B(Ts,Tt)=\lambda\mathcal B(s,t),\quad\lambda\ne0.
\]
With $L=\Lambda^2R$, one has
\begin{equation}
 A_{T\psi}=\lambda LA_\psi L^{-1},\qquad
 F_{16}(T\psi)=\lambda^8F_{16}(\psi),\qquad
 P_{18}(T\psi)=\lambda^9P_{18}(\psi).
 \label{rank:relative-invariant}
\end{equation}
The elementary units above, even particle--hole products, and $\Lambda g$
all satisfy these hypotheses, with their stated form multipliers.
\end{lemma}
\begin{proof}
The complete source map obeys $J_{T\psi}=TJ_\psi L^{-1}$.
The transformation of $H$ follows by applying the same form identity to
both of its arguments. Since $L$ preserves $G$, raising an index gives
the first formula. Hence $s_{T\psi}=\lambda s_\psi$ and
$t_k(T\psi)=\lambda^k t_k(\psi)$; substitution gives the other formulas.
For elementary units the Clifford conjugation is precisely the shear
\eqref{rank:elementary-shear}, whose graph is alternating, so it preserves
$q$ and $b$. Two-contraction units are dual. The CAR give the orthogonal
reflection, up to an overall minus sign, for $\tau_i$. Equation
\eqref{rank:gl-intertwining} gives the direction change
$(f,v)\mapsto(f\circ g^{-1},gv)$ for $\Lambda g$. The form multipliers
were established in \eqref{rank:bilinear-adjoint}--\eqref{rank:gl-multiplier}.
\end{proof}

\subsection{The two complete coordinate evaluations}

We give the coordinates explicitly so that the finite identities below
can be checked independently of a spectral or invariant-theoretic
classification. Write $u_i^+=(0,e_i)$ and $u_i^-=(f_i,0)$ in $U$, and
order the sixteen directions as
$z_0,\ldots,z_{15}=(u_0^+,\ldots,u_7^+,u_0^-,\ldots,u_7^-)$.
An integer $0\le m<256$ encodes an occupied subset through its bits $m_i$;
$e_m$ is the increasing exterior product on that subset. Put
\[
 \kappa(m)=(-1)^{\sum_{i=0}^7(i+1)m_i},\qquad
 X_i(m)=m\mathbin{\mathrm{xor}}2^{i\bmod8}.
\]
If $i<8$, let $\gamma_i(m)=(-1)^{\sum_{r<i}m_r}$ when $m_i=0$ and zero
otherwise. If $i\ge8$, use the same sign with $i$ replaced by $i-8$, and
require $m_{i-8}=1$. Thus $c(z_i)e_m=\gamma_i(m)e_{X_i(m)}$.
For a sorted pair $\alpha=(i,j)$ set
\[
 X_\alpha=X_iX_j,\qquad
 j_\alpha(m)=2\gamma_j(m)\gamma_i(X_j(m))
                   -\mathbf1_{j=i+8}.
\]
Then $2J_{e_m}(\alpha)=j_\alpha(m)e_{X_\alpha(m)}$.
Replace each direction in $\alpha$ by its dual, reorder increasingly,
and call the resulting pair $\alpha^\vee$ and the reordering sign
$\epsilon_\alpha$. The inverse of $G$ is exactly this signed dual-pair
operation. Therefore, for $\psi=\sum_{m=0}^{255}\psi_me_m$,
\begin{equation}
 (A_\psi)_{\alpha\beta}
 =\frac{\epsilon_\alpha}{4}\sum_{m=0}^{255}
 \psi_m j_{\alpha^\vee}(m)\kappa(X_{\alpha^\vee}(m))
 j_\beta(v)\psi_v,\qquad
 v=X_\beta\bigl(255\mathbin{\mathrm{xor}}X_{\alpha^\vee}(m)\bigr).
 \label{rank:entry-formula}
\end{equation}
This is just \eqref{rank:full-gram} expanded in the original occupation
basis. The factor $1/4$ accounts for the two occurrences of $2J$.
Every encoded subset and every direction in this formula stays in the original
finite domain.

\begin{lemma}\label{rank:normal-certificate}
Let
\[
 \Omega_*=(1+e_0e_1)(1+e_2e_3)(1+e_4e_5)(1+e_6e_7),
 \qquad \eta=\alpha1+\beta\mathrm{vol}+\gamma\Omega_*.
\]
For arbitrary $\alpha,\beta,\gamma\in\C$, set
$x=\alpha\beta$, $y=\alpha\gamma$, $z=\beta\gamma$,
$s=x+y+z$, and $w=xyz$. Then
\begin{align}
 s_\eta&=s,\label{rank:normal-halfpair}\\
 t_4(\eta)&=312s^4-1728sw,\label{rank:normal-t4}\\
 t_6(\eta)&=4152s^6-69984s^3w+36288w^2.
 \label{rank:normal-t6}
\end{align}
In particular $F_{16}(\eta)=P_{18}(\eta)=0$.
\end{lemma}
\begin{proof}
Here is a complete common basis for the matrix calculation. Partition the
mode indices into $B_r=\{2r,2r+1\}$, $0\le r<4$, and put
$\bar u=u\mathbin{\mathrm{xor}}1$, $\sigma_u=(-1)^u$.
There are eight zero directions
\[
 u_{2r}^+\wedge u_{2r+1}^-,\quad
 u_{2r+1}^+\wedge u_{2r}^-\qquad(0\le r<4).
\]
For each $r<s$, $i\in B_r$, $j\in B_s$, take the signed quartet
\begin{equation}
 \left(-u_i^+\wedge u_j^+,
 -\sigma_i u_j^+\wedge u_{\bar i}^-,
 \sigma_j u_i^+\wedge u_{\bar j}^-,
 -\sigma_i\sigma_j u_{\bar i}^-\wedge u_{\bar j}^-\right).
 \label{rank:quartet}
\end{equation}
These give 24 four-dimensional blocks. The remaining sixteen directions
are, for each $r$, the four entries
\begin{equation}
 \left(u_{2r}^+\wedge u_{2r+1}^+,\quad
 u_{2r}^+\wedge u_{2r}^-,\quad
 u_{2r+1}^+\wedge u_{2r+1}^-,\quad
 u_{2r}^-\wedge u_{2r+1}^-\right).
 \label{rank:central-basis}
\end{equation}
This list contains every pair of the sixteen split directions exactly once.
Substitution in \eqref{rank:entry-formula} gives the following whole matrix:
\begin{equation}
 A_\eta\simeq0_8\oplus T^{\oplus24}\oplus
       \bigl(T\otimes(I_4-P)+R\otimes P\bigr),\qquad
 P=\tfrac14\mathbf1\mathbf1^{\mathsf T},
 \label{rank:normal-blocks}
\end{equation}
where the last tensor factors index the four blocks in
\eqref{rank:central-basis}, and
\begin{equation}
 T=\begin{pmatrix}
 s&z&z&2z\\-y&0&0&-z\\-y&0&0&-z\\2y&y&y&s
 \end{pmatrix},\qquad
 R=\begin{pmatrix}
 s&-3z&-3z&-6z\\3y&2s&2s&3z\\3y&2s&2s&3z\\-6y&-3y&-3y&s
 \end{pmatrix}.
 \label{rank:TR}
\end{equation}
The square terms vanish for a structural reason. At the vacuum,
$J_1(W)\subseteq\Lambda^0E\oplus\Lambda^2E$, whose Chevalley pairing
with itself is zero in dimension eight. Hence $A_1=0$.
The even particle--hole product sends $1$ to $\mathrm{vol}$, and the
elementary creation product sends $1$ to $\Omega_*$.
Lemma~\ref{rank:covariance} gives $A_{\mathrm{vol}}=A_{\Omega_*}=0$.
Thus the coefficients of $\alpha^2,\beta^2,\gamma^2$ in
\eqref{rank:entry-formula} are zero. The coefficients of
$\alpha\beta,\alpha\gamma,\beta\gamma$ include both source orders.
The only nonzero source subsets are unions of the four pairs $B_r$.
On \eqref{rank:quartet}, evaluating these three coefficients yields $T$;
on the central directions $(i,r),(j,s)$ it yields
$T_{ij}\mathbf1_{r=s}+(R_{ij}-T_{ij})/4$; all other cross blocks vanish.
These rules, together with the eight zero directions, specify all
$120^2$ entries. Pair-block permutations reduce the quartet substitution
to the four parity choices for $i,j$; their signs are precisely those in
\eqref{rank:quartet}.

Since $P$ and $I-P$ are complementary projections of ranks one and three,
\eqref{rank:normal-blocks} gives, for every positive integer $k$,
\[
 t_k(\eta)=27\Tr(T^k)+\Tr(R^k).
\]
The trace computation can be performed with the following matrix
identities, obtained by multiplying \eqref{rank:TR}:
\begin{align*}
 T^4&=2sT^3-s^2T^2+4wT,\\
 R^4&=6sR^3-9s^2R^2+(4s^3-108w)R.
\end{align*}
The initial traces are
\[
 \begin{array}{c|ccc}
       &\Tr(\cdot)&\Tr((\cdot)^2)&\Tr((\cdot)^3)\\\hline
 T&2s&2s^2&2s^3+12w\\
 R&6s&18s^2&66s^3-324w
 \end{array}.
\]
Multiplying each recurrence by $I,T,T^2$ or $I,R,R^2$ and taking traces
gives \eqref{rank:normal-t4}--\eqref{rank:normal-t6}.
Equation \eqref{rank:chevalley-sign} gives
$\mathcal B(\eta,\eta)=2(\alpha\beta+\alpha\gamma+\beta\gamma)$.
Substituting these three expressions in \eqref{rank:polynomial} and
collecting powers of $s,w$ gives zero. All formulas are polynomial, so no
parameter is required to be nonzero.
\end{proof}

\begin{lemma}\label{rank:target-certificate}
For
\[
 v_A=e_0e_1e_2e_3,\quad v_B=e_4e_5e_6e_7,\quad
 \Psi=(1+v_A)(1+v_B),
\]
one has
\begin{equation}
 s_\Psi=2,\qquad t_4(\Psi)=3648,\qquad t_6(\Psi)=57408,
 \qquad P_{18}(\Psi)=18063360\ne0.
 \label{rank:target-values}
\end{equation}
\end{lemma}
\begin{proof}
Use \eqref{rank:entry-formula}, now with occupied-subset encodings $0,15,240,255$,
all with coefficient one. The 64 directions between the two
four-mode split spaces form an identity block. In each four-mode block,
the twelve directions $u_i^+\wedge u_j^-$ with $i\ne j$ give zero.
Pair each $u_i^+\wedge u_j^+$ with $u_k^-\wedge u_\ell^-$, where $k<\ell$ are
the complementary two indices of that four-mode block. Its matrix is
\[
 \begin{pmatrix}2&2(-1)^{i+j}\\2(-1)^{i+j}&2\end{pmatrix}.
\]
There are twelve such two-dimensional blocks in total. Finally the four
directions $u_i^+\wedge u_i^-$ in each block have the all-ones matrix of
size four. Entries between these blocks vanish by the same formula.
Thus the eigenvalues of $A_\Psi$ are $1$ with multiplicity $64$, $4$ with
multiplicity $14$, and $0$ with multiplicity $42$. Hence
$t_k(\Psi)=64+14\cdot4^k$ for $k>0$. The four complementary source pairs
give $\mathcal B(\Psi,\Psi)=4$. Substitution in
\eqref{rank:polynomial} proves \eqref{rank:target-values}.
\end{proof}

\subsection{All three summands, including the boundary}

The generic normal configuration in Lemma~\ref{rank:normal-certificate}
has the common symplectic symmetry studied by
Manivel~\cite[Proposition~1]{manivel2009}. To exclude an arbitrary
three-term decomposition, one must also treat third terms outside that
nondegenerate chart. The following construction does so on the original
vector by a polynomial specialization.

\begin{lemma}\label{rank:common-pair}
If $g_1,g_2$ are nonzero even pure spinors and
$\mathcal B(g_1,g_2)\ne0$, a common invertible transformation of the kinds
in Lemma~\ref{rank:covariance} sends them to
$\alpha1,\beta\mathrm{vol}$ with $\alpha\beta\ne0$.
The transformation acts simultaneously on any other summands.
\end{lemma}
\begin{proof}
Choose a nonzero coefficient $(g_1)_I$. The set $I$ is even.
The product $\tau_I$ toggles precisely the occupation set $I$, so the
vacuum coefficient of $\tau_Ig_1$ is $\pm(g_1)_I\ne0$.
Apply its inverse chart unit from Lemma~\ref{rank:chart} to the whole
configuration. The first term is now $\alpha1$.
The second term has top coefficient $\beta\ne0$, since the common
transformation has preserved $\mathcal B$.

Let $C=\tau_0\cdots\tau_7$. The CAR give $C^2=\Id$ and
$C1=\mathrm{vol}$, $C\mathrm{vol}=1$.
The vector $Cg_2'$ has vacuum coefficient $\beta$, so
$Cg_2'=\beta\mathcal U1$ for a chart unit $\mathcal U$.
Consequently $g_2'=\beta C\mathcal UC\,\mathrm{vol}$.
Conjugation by $C$ changes two-creation factors into two-contraction
factors. Their inverse product fixes the vacuum and takes the second
term to $\beta\mathrm{vol}$. Apply that same product to all terms.
\end{proof}

The first part of this proof also sends any even pure spinor to a
nonzero vacuum multiple. Its self-pairing is therefore zero, by the
nonzero form multiplier and $\mathcal B(1,1)=0$.

\begin{proposition}\label{rank:three-vanishing}
For every three nonzero even pure spinors $g_1,g_2,g_3\in\mathcal F_8^+$,
\[
 P_{18}(g_1+g_2+g_3)=0.
\]
\end{proposition}
\begin{proof}
Put $x=g_1+g_2+g_3$. If $s_x=0$, the assertion follows from
\eqref{rank:polynomial}. Otherwise, since all self-pairings vanish,
some pair has nonzero pairing. By Lemma~\ref{rank:common-pair} and
relative covariance it suffices to consider
\[
 x=\alpha1+\beta\mathrm{vol}+s,
\]
where $s$ is an arbitrary nonzero even pure spinor.

Choose a nonzero coefficient $s_I$, order $I$ increasingly, and divide it
into $k=|I|/2$ successive pairs $(i_r,j_r)$ with $i_r<j_r$.
Define the actual polynomial curve
\begin{equation}
 h(t)=\prod_{r=1}^k(\Id+t\iota_{f_{j_r}}\iota_{f_{i_r}})s,
 \qquad a(t)=[h(t)]_0.
 \label{rank:curve-first}
\end{equation}
The factors commute, square-zero generators give their inverses by
$t\mapsto-t$, and their Clifford shears preserve purity. Thus every
$h(t)$ is nonzero and pure. The coefficient of $t^k$ in $a(t)$ is
exactly $s_I$: all chosen contractions produce a vacuum only from the
complete original support $I$, and the successive ordered-pair convention
gives sign $+1$. Hence $a(t)$ is a nonzero polynomial. If $k=0$, this
statement simply reads $a(t)=s_\varnothing\ne0$.

Put $\omega_*=e_0e_1+e_2e_3+e_4e_5+e_6e_7$ and
\begin{equation}
 r(t,u)=\prod_{j=0}^3(\Id+u\varepsilon_{2j}\varepsilon_{2j+1})h(t).
 \label{rank:curve-second}
\end{equation}
This is again nonzero and pure for every $(t,u)$, and $r(0,0)=s$.
Its first two even components are
\[
 [r(t,u)]_0=a(t),\qquad [r(t,u)]_2=[h(t)]_2+u\,a(t)\omega_*.
\]
Let $K(\omega):E^*\to E$ be $f\mapsto\iota_f\omega$ and
$J=K(\omega_*)$, an invertible alternating matrix. The polynomial
\begin{equation}
 \Delta(t,u)=a(t)\det\bigl(K([h(t)]_2)+u\,a(t)J\bigr)
 \label{rank:discriminant}
\end{equation}
is nonzero: its coefficient of $u^8$ is $a(t)^9\det J$, nonzero in
$\C[t]$. At every point where $\Delta\ne0$, the vacuum coefficient is
nonzero and \eqref{rank:graph-coordinates} gives a nondegenerate graph
\[
 \mathcal D=-K\bigl([r(t,u)]_2/a(t)\bigr).
\]
The symplectic-basis construction after \eqref{rank:graph-gl} gives a
$g\in\GL(E)$ taking it to the graph of $\Omega_*$.
Lemma~\ref{rank:chart} and graph uniqueness then give
$\Lambda g\,r(t,u)=a(t)\Omega_*$.
The same $\Lambda g$ fixes $1$ and multiplies $\mathrm{vol}$ by $\det g$.
It therefore sends the entire sum
$\alpha1+\beta\mathrm{vol}+r(t,u)$ into the family of
Lemma~\ref{rank:normal-certificate}. Relative covariance proves
\[
 F_{16}\bigl(\alpha1+\beta\mathrm{vol}+r(t,u)\bigr)=0
 \quad\text{whenever }\Delta(t,u)\ne0.
\]

The left side is a polynomial $P(t,u)$. The product $P\Delta$ vanishes
at every point of $\C^2$, hence is the zero polynomial: a polynomial over
an infinite field is determined by all of its evaluations, by applying
the one-variable root bound successively to the two variables.
Since $\C[t,u]$ is a domain and $\Delta\ne0$, we have $P=0$.
Evaluating at $(0,0)$ recovers the exact original $s$, so
$F_{16}(\alpha1+\beta\mathrm{vol}+s)=0$.
There is no limiting assertion about the existence of a decomposition;
only the vanishing of one already defined polynomial is extended.
\end{proof}

\begin{corollary}\label{rank:two-copy}
The vector $\Psi=(1+v_A)(1+v_B)$ has Gaussian rank four.
\end{corollary}
\begin{proof}
Its four original occupation terms give the upper bound. A representation
with at most three Gaussian terms is a representation by even pure
spinors. If necessary, split one nonzero term into two nonzero scalar
multiples to obtain exactly three terms. Proposition~\ref{rank:three-vanishing}
would then give $P_{18}(\Psi)=0$, contrary to
Lemma~\ref{rank:target-certificate}.
\end{proof}

\begin{example}\label{rank:cross-pairing-example}
Each vector
\[
 \prod_{i=0}^3(1+z_i e_ie_{i+4}),\qquad z_i\in\C,
\]
is an eight-mode Gaussian with pairing across the two original four-mode
blocks. Such vectors are legitimate summands in a candidate shorter
expansion of $\Psi$. The preceding argument does not restrict them to
block products. It transports the whole proposed sum by one common map;
independently normalizing its three summands would destroy the sum identity
and would not imply the polynomial obstruction.
\end{example}

\subsection{Arbitrary four-mode factors}

\begin{lemma}\label{rank:four-mode-normal}
Every nonzero four-mode even vector is either Gaussian, of rank one, or
is carried to a nonzero scalar multiple of $1+\mathrm{vol}_4$ by an
invertible transformation preserving the Gaussian cone. In the latter
case its rank is two. The required transformations can be chosen as
complex Spin transformations followed by nonzero scalars.
\end{lemma}
\begin{proof}
An even particle--hole product makes a nonzero occupation coefficient a
nonzero vacuum coefficient $a$. Write the resulting vector as
\[
 v=a1+\omega+b\mathrm{vol}_4,
 \qquad \omega\in\Lambda^2\C^4.
\]
Choose the six elementary creation factors whose degree-two part is
$\omega/a$, and call their product $\mathcal U$.
There is a scalar $d$ such that
$\mathcal U1=1+\omega/a+d\mathrm{vol}_4$.
Positive-degree creators kill the volume, so
\[
 \mathcal U^{-1}v=a1+c\mathrm{vol}_4,\qquad c=b-ad.
\]
If $c=0$, the vector is Gaussian by \eqref{rank:cone}. If $c\ne0$, it is
not pure by Example~\ref{rank:four-mode-example}, and has rank two.
A diagonal $g\in\GL_4(\C)$ with $\det g=a/c$ takes it to
$a(1+\mathrm{vol}_4)$.

Here are explicit Spin lifts for all the operations, with split directions
$u_i^+,u_i^-$ defined as above in four modes. A particle--hole operator is
$c(u_i^++u_i^-)$, and $q(u_i^++u_i^-)=1$, so an even product belongs to
Spin. For an elementary creation unit choose $k$ distinct from $i,j$ and
put $v=u_k^++u_k^-$ and $w=v+u_j^+$. Both have quadratic norm one and
are orthogonal to $u_i^+$. In the Clifford algebra,
\[
 (v-zu_i^+)v\,(w+zu_i^+)w
 =(1-zu_i^+v)(1+zu_i^+w)=1+zu_i^+u_j^+.
\]
All four vector factors have quadratic norm one. This explicitly realizes
$1+z\varepsilon_i\varepsilon_j$ as a Spin action. For the last diagonal
scaling, choose $r\in\C^*$ with $r^2=a/c$. On one mode let
$P_0=\iota_i\varepsilon_i$ and $P_1=\varepsilon_i\iota_i$.
The two Clifford vectors $ru_i^++r^{-1}u_i^-$ and $u_i^++u_i^-$ have
quadratic norm one. Their product acts as $r^{-1}P_0+rP_1$ and
conjugates creation and contraction in that mode by factors
$r^2,r^{-2}$. Its spin action is $r^{-1}\Lambda g$ for the diagonal
$g$ used above. This supplies the claimed Spin transformation up to a
nonzero scalar without using a classification of half-spin orbits.
\end{proof}

\begin{proof}[Proof of Theorem~\ref{rank:main}]
First suppose both four-mode factors are non-Gaussian.
Lemma~\ref{rank:four-mode-normal} gives local complex Spin operations
and scalars taking them to $1+\mathrm{vol}_4$.
For the orthogonal direct sum $U_A\oplus U_B$, the natural map
\begin{equation}
 \Spin(U_A)\times\Spin(U_B)\longrightarrow\Spin(U_A\oplus U_B),
 \qquad (x,y)\longmapsto xy,
 \label{rank:local-spin}
\end{equation}
is a homomorphism because the even Clifford subalgebras commute.
It is not asserted to be injective: $(-1,-1)$ maps to the identity.
On the even--even Fock tensor factor its action is the usual tensor action.
Thus the local operations induce one invertible eight-mode transformation
preserving the entire physical Gaussian cone by \eqref{rank:cone}.
They take $\psi\otimes\phi$ to a nonzero multiple of $\Psi$.
Corollary~\ref{rank:two-copy} gives rank four, which is the product of
the two local ranks.

If one factor is Gaussian, use a local compact Gaussian operation and a
scalar to make it the vacuum. Tensoring a minimal decomposition of the
other factor with that vacuum gives one inequality. For the reverse
inequality, take any eight-mode Gaussian decomposition of
$1_A\otimes\phi$ and apply the same vacuum coefficient functional on all
four modes of $A$ to every summand. This is the composition of their
vacancy projections followed by removal of those vacuum modes.
The nonzero projected terms remain pure by the isotropic Clifford argument
following \eqref{rank:cone}. To see that removal preserves purity, write a
nonzero projected vector as $1_A\otimes t_B$. The $A$-degree-one and
$A$-degree-zero parts of its Clifford equation are independent, so
\[
 \Ann(1_A\otimes t_B)=E_A^*\oplus\Ann(t_B).
\]
Its dimension is eight, hence $\dim\Ann(t_B)=4$. Parity is even, and
\eqref{rank:cone} makes $t_B$ a physical Gaussian vector.
After discarding zero terms, the projected sum is a Gaussian decomposition
of the exact original $\phi$, with no more terms. This proves equality
of ranks. When both factors are Gaussian, their tensor product is Gaussian
by the compact version of \eqref{rank:local-spin}. All cases of
\eqref{rank:multiplicativity} follow.
\end{proof}

\begin{corollary}\label{rank:second-magic-state}
For the four-mode state
\[
 \widetilde M=\frac{\sqrt3\,|0000\rangle
   +\sqrt2\,|0011\rangle+\sqrt2\,|1100\rangle+|1111\rangle}{\sqrt8},
\]
one has $\chi_G(\widetilde M)=2$ and
$\chi_G(\widetilde M^{\otimes2})=4$.
\end{corollary}
\begin{proof}
Before normalization its degree-two part is the sum of two disjoint
pair terms, each with coefficient $\sqrt2$. Their product has top
coefficient $2$, whereas the product of its vacuum and top coefficients
is $\sqrt3$. Equivalently, the reduction in
Lemma~\ref{rank:four-mode-normal} gives
$c=1-2/\sqrt3\ne0$. The local rank is therefore two, and
Theorem~\ref{rank:main} gives the product rank.
\end{proof}
This state also occurs in the approximate-decomposition experiments of
\cite{cudby2023}. The corollary concerns exact rank; it does not determine
the minimum number of Gaussian terms at a positive approximation error.

The current Lean development verifies the original rank bounds
$2\le\chi_G(\Psi)\le4$ and the complete finite-coordinate identities above;
it does not yet formalize the rank-four theorem end to end.
The precise verified statements and remaining identifications are described
in Section~\ref{sec:formalization}.

\section{Entropy deficits and coherent measurement capacity}
\label{app:section}

The distribution comparison also gives quantitative information about
the input state and the measurement it defines.  The first application
relates the Wehrl entropy deficit to fractional occupation numbers.
The second determines the mutual-information capacity of the coherent
measurement, including repeated measurements on entangled inputs.

\subsection{An occupation-number bound for the entropy deficit}

\begin{proposition}
\label{app:deficit}
Let $2\le p\le n$, let $\psi\in\Lambda^p\C^n$ be a unit vector,
write $\rho=|\psi\rangle\langle\psi|$, and put $D=\binom np$.
With the one-particle matrix $\gamma_\rho$ from
\eqref{ext:one-body}, and the coherent moments and entropy constants
of Section~\ref{ent:section},
\begin{equation}
 S_W(\rho)-s_{p,n}
 \ge
 \frac{D\,m_{p-1,n-1}(2)}{2n(n+1)}
 \Tr\bigl[\gamma_\rho(\Id-\gamma_\rho)\bigr].
 \label{app:entropy-deficit}
\end{equation}
\end{proposition}
\begin{proof}
The original occupation conditioning \eqref{ext:conditioning},
followed by the lower-sector comparison with the test $t^2$, gives
\begin{equation}
 \int Q_\rho^2\,d\nu_{p,n}
 \le m_{p-1,n-1}(2)\,
       \E_u\langle u,\gamma_\rho u\rangle^2,
 \label{app:quadratic-conditioning}
\end{equation}
where $u$ is uniform on the complex unit sphere.  Zero conditional
branches contribute zero, so this formula also covers singular
$\rho$ and vanishing occupation probabilities.

For every Hermitian matrix $A$, diagonalization and unitary invariance
give
\begin{equation}
 \E_u\langle u,Au\rangle^2
 =\frac{(\Tr A)^2+\Tr(A^2)}{n(n+1)}.
 \label{app:sphere-second-moment}
\end{equation}
Indeed, normalize a vector of independent standard complex Gaussian
coordinates.  The squared moduli of its normalized coordinates are
uniform on the probability simplex.  Integrating its monomials gives
$\E|u_i|^4=2/[n(n+1)]$ and
$\E|u_i|^2|u_j|^2=1/[n(n+1)]$ for $i\ne j$.
Expanding the square in an eigenbasis proves
\eqref{app:sphere-second-moment}; the formula also holds for $n=1$.

A coherent Slater input has one-particle matrix $\pi_W$, with
$\Tr\pi_W=\Tr(\pi_W^2)=p$.  Its conditional states are Slater,
so \eqref{app:quadratic-conditioning} is an equality for that input.
Since $\Tr\gamma_\rho=p$, subtraction yields
\begin{equation}
 m_{p,n}(2)-\int Q_\rho^2\,d\nu_{p,n}
 \ge\frac{m_{p-1,n-1}(2)}{n(n+1)}
          \Tr\bigl[\gamma_\rho(\Id-\gamma_\rho)\bigr].
 \label{app:moment-deficit}
\end{equation}
Finally, $t\log t-\tfrac12t^2$, extended continuously at zero,
is convex on $[0,1]$, since its second derivative on $(0,1]$ is
$t^{-1}-1\ge0$.  Theorem~\ref{ext:convex-order} applied to this
function gives
\[
 S_W(\rho)-s_{p,n}
 \ge\frac D2\left(m_{p,n}(2)-\int Q_\rho^2\,d\nu_{p,n}\right).
\]
Combining this with \eqref{app:moment-deficit} proves the result.
\end{proof}

The right-hand side sums $\lambda_j(1-\lambda_j)$ over the natural
occupation numbers.  It vanishes exactly when the one-particle matrix
is a projection, which by the equality argument of
Theorem~\ref{ext:convex-order} forces a pure Slater state.
The coefficient in \eqref{app:entropy-deficit} is not asserted to be
optimal.  For $p=n$ the sector is one-dimensional and both sides
are zero.

\subsection{The information capacity of the coherent POVM}

Fix one of the coherent orbits considered above.  Let $\mathcal H$
be its representation space, $D=\dim\mathcal H$, $\nu$ its invariant
probability, and $s_*$ its coherent Wehrl entropy from
Theorem~\ref{ent:minima}.  The resolution of the identity defines
the positive operator-valued measure
\begin{equation}
 E(dx)=D|\Omega_x\rangle\langle\Omega_x|\,d\nu(x).
 \label{app:povm}
\end{equation}
An input density operator $\rho$ produces the density
$f_\rho(x)=D Q_\rho(x)$ relative to $\nu$.
The \emph{measurement capacity} is the supremum of $I(Z:X)$ over
measurable input ensembles $z\mapsto\rho_z$ with a probability law for the
classical label $Z$; the measurement itself remains fixed.

\begin{proposition}[Capacity and repeated inputs]
\label{app:capacity}
For the POVM \eqref{app:povm},
\begin{equation}
 C_{\mathrm{meas}}=\log D-s_*.
 \label{app:measurement-capacity}
\end{equation}
For every positive integer $\ell$, its $\ell$-fold product measurement
has capacity $\ell(\log D-s_*)$, allowing arbitrary joint input
ensembles of density operators on $\mathcal H^{\otimes\ell}$.
In particular, this assertion includes entangled inputs.
\end{proposition}
\begin{proof}
For a probability density $f$ with respect to the probability measure
$\nu$, write $h_\nu(f)=-\int f\log f\,d\nu$.
The normalization of Wehrl entropy gives
\begin{equation}
 h_\nu(f_\rho)=S_W(\rho)-\log D\ge s_*-\log D.
 \label{app:output-entropy}
\end{equation}
Also $h_\nu(f)\le0$, by convexity of $t\log t$ and $\int f\,d\nu=1$.
For an ensemble $(\rho_z,\pi(dz))$, put
$\overline f=\int f_{\rho_z}\,\pi(dz)$.  Then
\[
 I(Z:X)
 =h_\nu(\overline f)-\int h_\nu(f_{\rho_z})\,\pi(dz)
 \le\log D-s_*.
\]
These entropy expressions are finite because all measurement output
densities are bounded by $D$.  Equality is attained by the invariant
ensemble of pure coherent projectors: its average input is $\Id/D$,
its average output density is one, and each input has Wehrl entropy
$s_*$.  A finite ensemble can attain the same value.  Indeed,
$\Id/D$ belongs to the convex hull of the compact orbit of coherent
projectors, and finite-dimensional convexity gives a finite convex
combination with this average; every member still has the same
minimum output entropy.

For repeated measurements, fix any density operator $\rho$ on
$\mathcal H^{\otimes\ell}$.  Its joint output density is
\[
 f_\rho(x_1,\ldots,x_\ell)
 =D^\ell\Tr\!\left[\rho\bigotimes_{j=1}^\ell
                |\Omega_{x_j}\rangle\langle\Omega_{x_j}|\right].
\]
For a history $(x_1,\ldots,x_{j-1})$ of positive marginal density,
define
\[
 K=\bigotimes_{i<j}\sqrt D\,
       |\Omega_{x_i}\rangle\langle\Omega_{x_i}|,
 \qquad
 \tau=
 \frac{\Tr_{1,\ldots,j-1}[(K\otimes\Id)\rho(K^*\otimes\Id)]}
      {\Tr[(K\otimes\Id)\rho(K^*\otimes\Id)]}.
\]
The operator $\tau$ is a density operator on the remaining factors.
Its first reduced density operator, say $\tau_j$, produces the
conditional density of $x_j$ by the same POVM:
$f(x_j\mid x_1,\ldots,x_{j-1})=D Q_{\tau_j}(x_j)$.
This follows by inserting the resolution of the identity in the
unmeasured factors and taking the indicated partial trace.
Histories with zero marginal density make no contribution.

Each conditional density is bounded by $D$ and satisfies
\eqref{app:output-entropy}.  The joint density is bounded by $D^\ell$.
Thus $f\log f$ is integrable for each joint and conditional law;
in particular, its negative part is bounded in integral by $1/e$
and its positive part by the corresponding logarithmic bound.
The entropy chain rule is therefore valid and gives
\begin{equation}
 h_{\nu^{\otimes\ell}}(f_\rho)
 =\sum_{j=1}^\ell
   \E\,h_\nu\bigl(f(\,\cdot\mid X_1,\ldots,X_{j-1})\bigr)
 \ge\ell(s_*-\log D).
 \label{app:repeated-entropy}
\end{equation}
Product coherent inputs attain this bound.  For an arbitrary joint
input ensemble, the entropy of its average output relative to
$\nu^{\otimes\ell}$ is at most zero.  The mutual-information identity
therefore gives the upper bound $\ell(\log D-s_*)$.
Independent copies of an optimal coherent ensemble attain it,
which proves the capacity formula even with joint encoding.
\end{proof}

\section{Formal verification}
\label{sec:formalization}

The Lean development accompanies the proofs above and records which parts
have been checked on their original mathematical objects. Its source and
verification evidence are available in the repository
\cite{repository2026}. The verified source snapshot is
\href{https://github.com/veridiscoverLab/fermionic-coherent-lean/tree/3164a3243fcc6cd675fa8a7493282a1f650a16d8}{\texttt{3164a3}};
the English manuscript was added subsequently without changing those proof
sources.

\subsection{End-to-end exterior-power theorems}

The formalization of Theorem~\ref{ext:convex-order} starts with
$\Lambda^p\C^n$, its actual unitary action, density operators, and normalized
Haar measure. It constructs creation and contraction, the conditional
density on the orthogonal complement, the double-Haar identity, and the
common occupation mixture. The induction and the strict equality argument
are proved from these constructions. In particular, the statement covers
every $0\le p\le n$, all mixed inputs, and every continuous convex test on
$[0,1]$; a strictly convex test has equality exactly at pure Slater
densities. No conditional-measure identity or full-density rigidity
statement is assumed as a separate hypothesis.

The formal Slater measure is the pushforward to the actual orbit of
rank-one Slater projectors. Its invariance, independence of the reference,
and agreement with group integration are proved. The entropy results then
use this same measurement. They include the sector dimension $\binom np$,
the normalized first moment, all positive real moments, the Wehrl minimum,
and the R\'enyi--Wehrl minimum for every real $q>0$, $q\ne1$, together with
the same equality characterization. These minimizer theorems are end to
end. Their coherent comparison values are defined by integrals; the
Beta-product evaluations, explicit Gamma products, harmonic-number
constants, and the limit as $q\to1$ have not been formalized.

The principal declarations are listed below. Each namespace has the prefix
\texttt{Fermionic.}; the repository prints the full types of all declarations,
including their quantifiers and domain assumptions.
\begin{center}
\small
\begin{tabular}{@{}p{0.28\linewidth}p{0.66\linewidth}@{}}
\toprule
Namespace & Main declarations \\
\midrule
\texttt{HusimiInduction} &
\texttt{exterior\_husimi\_convex\_order}\newline
\texttt{exterior\_husimi\_equality\_iff} \\[3pt]
\texttt{SlaterMeasure} &
\texttt{orbit\_husimi\_convex\_order}\newline
\texttt{orbit\_husimi\_equality\_iff} \\[3pt]
\texttt{ExteriorEntropy} &
\texttt{wehrl\_minimum}, \texttt{wehrl\_eq\_iff\_slater}\newline
\texttt{renyiWehrl\_minimum}\newline
\texttt{renyiWehrl\_eq\_iff\_slater} \\
\bottomrule
\end{tabular}
\end{center}

\subsection{The Gaussian-rank boundary}

For the actual even Spin-vacuum-orbit Gaussian dictionary on eight modes,
the formalized results give, for the original occupation vector $\Psi$
defined in Section~\ref{sec:gaussian-rank},
\[
 2\le\chi_G(\Psi)\le4.
\]
They include the actual Clifford representation, the four-term Gaussian
decomposition, purity of Gaussian orbit vectors, the vacuum-chart
annihilator calculation, and nonpurity of the target. Formal scale
invariance gives the same interval for $\Psi/2$. Its occupation-space
identification with $M^{\otimes2}$ is made in the paper; a separate
Hilbert tensor-product identification is not part of the verified roots.
Separately, the full
normal-family certificate in Section~\ref{sec:gaussian-rank} is checked for
arbitrary complex parameters, including zero parameters. The encoded
target has the nonzero certificate value $18063360$.

These results do not yet constitute an end-to-end formalization of
$\chi_G(M^{\otimes2})=4$. The remaining Lean work is the original
Fock-to-coordinate identification, covariance of the certificate under the
actual common transformations, and the reduction of every three-term
pure-spinor sum, including its degenerate cases. Section~\ref{sec:gaussian-rank}
gives the mathematical argument for these steps. The half-spin and odd-spin
convex-order theorems, their entropy consequences, and the applications in
Section~\ref{app:section} likewise have paper proofs here, rather than
completed end-to-end Lean proofs. The repository's
\href{https://github.com/veridiscoverLab/fermionic-coherent-lean/blob/main/TODO.md}{\texttt{TODO.md}}
records the outstanding formal statements and their paper-level derivations.

\subsection{Verification procedure and foundations}

The recorded runs use Lean~4.29.0 and mathlib at revision
\texttt{8a178386ffc0f5fef0b77738bb5449d50efeea95}.
All Lean and Lake execution took place on a CAB server through SSH.
Three suites passed clean project compilation with warnings treated as
errors, inspection of all selected declaration types and their transitive
axioms, and a fresh kernel replay of the imported proof environment.
The exterior suite contains 20 modules and 309 theorems; the rank
foundation and certificate suite contains 33 modules and 257 theorems;
the vacuum-chart suite contains 12 modules and 150 theorems. These suites
overlap and were verified in three separate runs. Their shared source
files are identical.

The proofs use no \texttt{sorry}, \texttt{admit}, added mathematical axiom,
or \texttt{native\_decide}. Finite certificate leaves use ordinary
kernel-checked decision procedures, and parameter identities use
kernel-checked algebraic proofs. The disclosed logical axioms are
\texttt{propext}, \texttt{Classical.choice}, and \texttt{Quot.sound}.
Fresh replay starts from an empty environment and rechecks the replayable
safe, non-partial constants with the same Lean kernel; axioms remain axioms.
It is not an independent implementation of the kernel. The dependency
versions are pinned, but their source files were not all recompiled from
scratch in each run. Source hashes, complete logs, and reproduction
instructions are included in the repository.

\begingroup
\small
\begin{thebibliography}{99}
\setlength{\itemsep}{3pt}
\setlength{\parsep}{0pt}
\raggedright

\bibitem{gnutzmann2001}
S.~Gnutzmann and K.~\.{Z}yczkowski,
R\'enyi--Wehrl entropies as measures of localization in phase space,
\emph{J. Phys. A: Math. Gen.} \textbf{34} (2001), 10123--10139.
\href{https://doi.org/10.1088/0305-4470/34/47/317}{doi:10.1088/0305-4470/34/47/317}.

\bibitem{sugita2002}
A.~Sugita,
Proof of the generalized Lieb--Wehrl conjecture for integer indices larger than one,
\emph{J. Phys. A: Math. Gen.} \textbf{35} (2002), L621--L626.
\href{https://doi.org/10.1088/0305-4470/35/42/105}{doi:10.1088/0305-4470/35/42/105}.
\href{https://arxiv.org/abs/nlin/0208007}{arXiv:nlin/0208007}.

\bibitem{sugita2003}
A.~Sugita,
Moments of generalized Husimi distributions and complexity of many-body quantum states,
\emph{J. Phys. A: Math. Gen.} \textbf{36} (2003), 9081--9103.
\href{https://doi.org/10.1088/0305-4470/36/34/310}{doi:10.1088/0305-4470/36/34/310}.

\bibitem{liebsolovej2016}
E.~H.~Lieb and J.~P.~Solovej,
Proof of the Wehrl-type entropy conjecture for symmetric $\mathrm{SU}(N)$ coherent states,
\emph{Commun. Math. Phys.} \textbf{348} (2016), 567--578.
\href{https://doi.org/10.1007/s00220-016-2596-9}{doi:10.1007/s00220-016-2596-9}.
\href{https://arxiv.org/abs/1506.07633}{arXiv:1506.07633}.

\bibitem{cudby2023}
J.~Cudby and S.~Strelchuk,
Gaussian decomposition of magic states for matchgate computations,
preprint, 2023; version 4, 2025.
\href{https://arxiv.org/abs/2307.12654v4}{arXiv:2307.12654v4}.

\bibitem{diaskoenig2024}
B.~Dias and R.~Koenig,
Classical simulation of non-Gaussian fermionic circuits,
\emph{Quantum} \textbf{8} (2024), 1350.
\href{https://doi.org/10.22331/q-2024-05-21-1350}{doi:10.22331/q-2024-05-21-1350}.
\href{https://arxiv.org/abs/2307.12912}{arXiv:2307.12912}.

\bibitem{reardon2024}
O.~Reardon-Smith,
The fermionic linear optical extent is multiplicative for 4 qubit parity eigenstates,
preprint, 2024.
\href{https://arxiv.org/abs/2407.20934}{arXiv:2407.20934}.

\bibitem{tarabunga2026}
P.~S.~Tarabunga,
Fermionic non-Gaussianity via Bell sampling: monotones and efficient quantum algorithms,
preprint, 2026; version 2.
\href{https://arxiv.org/abs/2606.05066v2}{arXiv:2606.05066v2}.

\bibitem{hackl2021}
L.~Hackl and E.~Bianchi,
Bosonic and fermionic Gaussian states from K\"ahler structures,
\emph{SciPost Phys. Core} \textbf{4} (2021), 025.
\href{https://doi.org/10.21468/SciPostPhysCore.4.3.025}{doi:10.21468/SciPostPhysCore.4.3.025}.
\href{https://arxiv.org/abs/2010.15518}{arXiv:2010.15518}.

\bibitem{manivel2009}
L.~Manivel,
On spinor varieties and their secants,
\emph{SIGMA} \textbf{5} (2009), 078, 22 pp.
\href{https://doi.org/10.3842/SIGMA.2009.078}{doi:10.3842/SIGMA.2009.078}.

\bibitem{repository2026}
veridiscoverLab,
Fermionic coherent states in Lean,
source code and verification evidence, 2026.
\url{https://github.com/veridiscoverLab/fermionic-coherent-lean}.

\end{thebibliography}
\endgroup

\end{document}
